\documentclass[11pt,letterpaper]{article}

\usepackage[margin=1in]{geometry}
\usepackage{amsmath,amssymb,amscd,amsthm,mathtools}
\usepackage{authblk}
\usepackage{booktabs}
\usepackage[hidelinks]{hyperref}
\usepackage{microtype}
\usepackage{enumitem}

\numberwithin{equation}{section}

\theoremstyle{plain}
\newtheorem{theorem}{Theorem}[section]
\newtheorem{lemma}[theorem]{Lemma}
\newtheorem{proposition}[theorem]{Proposition}
\newtheorem{corollary}[theorem]{Corollary}
\theoremstyle{definition}
\newtheorem{definition}[theorem]{Definition}
\theoremstyle{remark}
\newtheorem{remark}[theorem]{Remark}

\newcommand{\NN}{\mathbb{N}}
\newcommand{\ZZ}{\mathbb{Z}}
\newcommand{\Xfin}{X_\infty^{\mathrm{fin}}}
\newcommand{\Xinf}{X_\infty}
\DeclareMathOperator{\supp}{supp}
\DeclareMathOperator{\Val}{Val}
\DeclareMathOperator{\NF}{NF}
\DeclareMathOperator{\Fold}{Fold}

\renewcommand\Authfont{\normalsize}
\renewcommand\Affilfont{\small}
\setlength{\affilsep}{0.4em}

\title{Finite-window rigidity in Fibonacci numeration}
\author[1]{Haobo Ma}
\author[2]{Wenlin Zhang\thanks{Corresponding author.}}
\affil[1]{AELF PTE LTD., \#14-02, Marina Bay Financial Centre Tower 1, 8 Marina Blvd, Singapore 018981, Singapore\\\texttt{auric@aelf.io}}
\affil[2]{National University of Singapore, Singapore\\\texttt{e1327962@u.nus.edu}}
\date{}

\begin{document}
\maketitle

\begin{abstract}
For each $m\ge 1$, let $\Fold_m:\{0,1\}^m\to X_m$ send a binary word to the
first $m$ digits of the Zeckendorf expansion of its Fibonacci value.  We show
that $\Fold_m$ is the normal-form map of a finite terminating confluent rewrite
system and that the family $(\Fold_m)_{m\ge 1}$ is compatible with restriction
in $m$.  For $m\ge 3$ we prove that the associated sliding block code $\Phi_m$
has the following exact reconstruction property: every block of $n\ge m$
consecutive output symbols has a unique binary preimage of length $n+m-1$.
Equivalently,
the block map $B_{m,n}:\{0,1\}^{n+m-1}\to\mathcal L_n(Y_m)$ is a bijection and
$\lvert\mathcal L_n(Y_m)\rvert=2^{n+m-1}$.
Equivalently, for $m\ge 3$ the residue vector modulo $F_{m+2}$ of the
overlapping local Fibonacci sums determines the entire binary block uniquely.

This block bijection gives local inverse rules with memory $r$ and anticipation
$m-1-r$ for every $0\le r\le m-1$, and the one-sided inverse memory $m-1$ is
best possible.  It also yields explicit formulas for Bernoulli measures pushed
forward from the full two-shift: if $\mu_{m,p}=(\Phi_m)_*\nu_p$, then
$\mu_{m,p}([a])=p^{|\omega(a)|_1}(1-p)^{n+m-1-|\omega(a)|_1}$ for
$a\in\mathcal L_n(Y_m)$, and
$H_{\mu_{m,p}}(Y_0^{n-1})=(n+m-1)h_2(p)$.  For Bernoulli input bits, this
arithmetic residue vector has an explicit finite-dimensional law and forms an
exact $m$-step Markov process.  We also show that the right Fischer cover of
$Y_m$ over the natural alphabet $X_m$ is the explicit suffix graph $G_m$, which
has $2^{m-1}$ states and $2^m$ edges.
\end{abstract}

\medskip
\noindent\textbf{Keywords.}
Zeckendorf expansion, Fibonacci numeration, finite-window rigidity, Bernoulli
transport, exact decoding, Markov order, finite-state normalization.

\medskip
\noindent\textbf{2020 Mathematics Subject Classification.}
Primary 11A63; Secondary 11B39, 68Q45, 37B10.


\section{Introduction}

Zeckendorf expansion singles out a canonical Fibonacci representation of every
positive integer.  The problem considered here is a finite-window version of
that principle.  For each $m\ge 1$ we consider the binary window space
\[
\Omega_m:=\{0,1\}^m.
\]
Let
\[
X_m:=\{w\in\{0,1\}^m:\text{$w$ contains no subword }11\},
\]
and define
\[
\Fold_m(\omega):=\pi_m(Z(N(\omega))),
\]
where $N(\omega)$ is the Fibonacci value of $\omega$, $Z(N(\omega))$ is its
Zeckendorf expansion, and $\pi_m$ records the first $m$ digits.  Thus
$\Fold_m$ replaces an arbitrary binary window by its canonical admissible
representative in Fibonacci numeration.

Two questions motivate the paper.  The first is local: can $\Fold_m$ be
described by a finite normalization procedure, and how does it behave when the
window length changes?  The second is reconstructive: if one slides this
normalization along a bi-infinite binary sequence, how much of the original
sequence can be recovered from the resulting list of overlapping folded
windows?  In number-theoretic terms, the second question asks whether the
vector of overlapping local Fibonacci weighted sums modulo $F_{m+2}$
determines an arbitrary binary block uniquely, and if so, what the sharp
observation length is.

We answer the first question by showing that $\Fold_m$ is the normal-form map
of a finite terminating confluent rewrite system; see
Proposition~\ref{prop:rewrite-newman} and
Corollary~\ref{cor:fold-order-independent}.  We also prove that the maps
$\Fold_m$ are compatible with restriction in the window length.  In particular,
the finite admissible languages form an inverse system whose limit is the
one-sided golden-mean shift; see Theorem~\ref{thm:inverse-limit-golden}.

To study the second question, for each $m$ we consider the sliding block code
\[
\Phi_m:\{0,1\}^{\ZZ}\to X_m^{\ZZ}
\]
with image $Y_m$, the language of overlapping normalized windows.  The main
result is that for every $m\ge 3$ and every $n\ge m$, each admissible block of
$n$ consecutive folded symbols has a unique raw lift of length $n+m-1$.  Thus
the threshold $m$ is exactly the point at which local normalization becomes
lossless on finite blocks.  This yields the block count
\[
\lvert\mathcal L_n(Y_m)\rvert=2^{n+m-1}
\qquad (n\ge m),
\]
and gives a family of local inverse rules obtained by reading any prescribed
coordinate of the unique length-$m$ lift.  The same block bijection also
describes the Bernoulli measures pushed forward from the full two-shift; see
Theorem~\ref{thm:block-bijection-threshold},
Corollary~\ref{cor:block-counts-threshold},
Corollary~\ref{cor:higher-block-identification}, and
Theorems~\ref{thm:Phi-m-conjugacy-threshold},
\ref{thm:bernoulli-transport-cylinder-law}, and
\ref{thm:exact-markov-order-bernoulli}.  At the one-sided endpoint this
recovers an inverse with memory $m-1$, and no smaller one-sided memory
suffices; see Corollary~\ref{cor:Phi-m-memory-threshold-exact}.

Thus the same integer $m$ governs several different aspects of the problem.  It
is the first scale at which overlapping folded windows determine the
underlying raw block, the Markov memory of the push-forward of Bernoulli
measures, the minimal synchronizing length of the natural presentation, and
the size parameter of the right Fischer cover.  The symbolic-dynamical
language serves mainly to record these consequences efficiently.

\subsection*{Related work}

The arithmetic background begins with Zeckendorf's theorem
\cite{Zeckendorf1972}.  Finite-state aspects of numeration systems were
developed systematically by Frougny and collaborators
\cite{Frougny1991FibonacciFiniteAutomata,Frougny1992NumbersAutomata,
Frougny1999OnlineAddition,FrougnySolomyak1992}, while algorithmic forms of
Zeckendorf arithmetic, including carry propagation and parallel normalization,
were studied in \cite{AhlbachUsatineFrougnyPippenger2013}.  For broader
numeration settings one may compare Ostrowski systems and their
automata-theoretic realizations
\cite{Berthe2001Ostrowski,HieronymiTerry2018OstrowskiAddition,
BaranwalaSchaefferShallit2021OstrowskiAutomatic}, as well as more general work
on non-standard and linear-recurrence numeration systems with finite-memory
automaton structure and normalization algorithms
\cite{AlloucheShallit2003AutomaticSequences,
FrougnyPelantovaSvobodova2013MinimalDigitSets,
Frougny2012NumbersAsStreamsLeiden,LegerskySvobodova2018EWM}.  Recent work has
also connected Zeckendorf representations with sequence-level statistical and
dynamical phenomena; see \cite{ManiboMiroRustTadeo2020ZeckendorfMixing}.

The distributional and combinatorial structure of Zeckendorf decompositions has
been investigated by Grabner, Tichy, and collaborators from the viewpoint of
digital problems and asymptotic distribution of digit sums in non-standard
numeration systems \cite{GrabnerTichy1995FibonacciDigitalSums}; by Miller and
collaborators, who established Gaussian behavior for the number of summands in
Zeckendorf representations \cite{MillerWang2012GaussianZeckendorf}; and from the
ergodic-theoretic perspective connecting beta-expansions and Parry conditions to
symbolic dynamics of numeration \cite{Parry1960BetaExpansions,
Renyi1957Representations}.

These strands of literature address different questions: recognition of legal
expansions, arithmetic operations in non-standard numeration systems,
automaticity, distributional behavior, or digit-sum statistics.  The present
paper does not study addition, parallel arithmetic, or asymptotic digit
statistics.  Instead it fixes a window length $m$, normalizes every window
locally in Fibonacci numeration, and asks when the original binary block can be
reconstructed exactly from the resulting local Zeckendorf data.  Equivalently,
the maps $\Phi_m$ form a family of finite-window factor maps attached to
Fibonacci numeration, and the later finite-state constructions record the
corresponding determinization data.

More concretely, the work of Frougny and collaborators shows how finite automata
recognize admissible expansions and implement normalization or arithmetic in
Fibonacci and related numeration systems.  Our question is different.  We begin
with an arbitrary binary block, apply truncated Zeckendorf normalization on
every overlapping window of length $m$, and ask when those local normal forms
still determine the original block exactly.  Unlike the above literature, the
present paper proves a sharp local congruence rigidity threshold and derives its
exact consequences for representation counting and random local residue
distributions.  We are not aware of an earlier result, even in the Fibonacci
setting, that identifies this sharp observation length for the recoverability
problem or the corresponding threshold in the natural deterministic
presentation.

This threshold is of arithmetic interest because it measures how much
information survives finite truncation of Zeckendorf normalization.  The answer
is not merely qualitative.  Once the window length reaches $m$, every
admissible folded block of length $n\ge m$ corresponds to exactly one binary
block of length $n+m-1$, so
$\lvert\mathcal L_n(Y_m)\rvert=2^{n+m-1}$.  The same bijection gives explicit
push-forward formulas for Bernoulli measures and an exact $m$-step Markov law
on the folded alphabet.

In purely number-theoretic terms, the main results can be stated as follows.
For each binary block $\omega=\omega_1\cdots\omega_L\in\{0,1\}^L$ with
$L=n+m-1$, the $t$-th local Fibonacci sum is
$S_t^{(m)}(\omega):=\sum_{j=0}^{m-1}\omega_{t+j}F_{j+2}$, and its residue
modulo $F_{m+2}$ is
$R_t^{(m)}(\omega):=S_t^{(m)}(\omega)\bmod F_{m+2}$.  The residue vector
$\mathcal R_{m,n}(\omega):=(R_1^{(m)}(\omega),\ldots,R_n^{(m)}(\omega))$
determines $\omega$ uniquely.  The symbolic-dynamical language provides a
convenient framework for organizing the proofs of this and the following
consequences.

\begin{theorem}[Main results]
\label{thm:intro-main}
For every $m\ge 3$, the following statements hold.
\begin{enumerate}[label=\textup{(\roman*)}]
\item\label{item:main-congruence-rigidity}
\textup{(Local Fibonacci congruence rigidity.)}
For every $n\ge m$ and $L=n+m-1$, the map
$\mathcal R_{m,n}:\{0,1\}^L\to\{0,1,\ldots,F_{m+2}-1\}^n$ that sends a binary
block to its vector of overlapping local Fibonacci sums modulo $F_{m+2}$ is
injective.  Equivalently, the congruence system
\[
\sum_{j=0}^{m-1}\omega_{t+j}F_{j+2}\equiv r_t\pmod{F_{m+2}}
\qquad (1\le t\le n)
\]
has at most one solution $\omega\in\{0,1\}^L$ for each
$(r_1,\ldots,r_n)\in\{0,1,\ldots,F_{m+2}-1\}^n$.
\item\label{item:main-block-bijection}
\textup{(Exact block reconstruction.)}
For every $n\ge m$, every admissible block of $n$ consecutive folded
symbols has a unique raw lift of length $n+m-1$.  Equivalently, the block map
$B_{m,n}:\{0,1\}^{n+m-1}\to\mathcal L_n(Y_m)$ is a bijection.
Consequently
$\lvert\mathcal L_n(Y_m)\rvert=2^{n+m-1}$ for every $n\ge m$, and the
$n$-block presentation of the normalized language is canonically the
$(n+m-1)$-block presentation of the full binary shift.
\item\label{item:main-bit-recovery}
\textup{(Exact bit recovery from residue windows.)}
Write $\mathfrak{R}_{m,m}$ for the set of realizable length-$m$ residue
vectors.  For each $0\le s\le m-1$, there exists a function
$\rho_{m,s}:\mathfrak{R}_{m,m}\to\{0,1\}$ such that
\[
\omega_{t+s}=\rho_{m,s}\bigl(R_t^{(m)}(\omega),\ldots,R_{t+m-1}^{(m)}(\omega)\bigr)
\]
for every $\omega\in\{0,1\}^{n+m-1}$ and every valid index $t$.  The minimum
number of consecutive residues required for any one-sided recovery rule is
exactly $m$.
\item\label{item:main-bernoulli}
\textup{(Explicit law and exact Markov order.)}
For every $p\in(0,1)$, if $\nu_p$ denotes the Bernoulli product law on
$\{0,1\}^{\ZZ}$ and $\mu_{m,p}:=(\Phi_m)_*\nu_p$, then every cylinder
$[a]_0\subset Y_m$ with $a\in\mathcal L_n(Y_m)$ and $n\ge m$ satisfies the
explicit formula
\[
\mu_{m,p}([a]_0)=p^{|\omega(a)|_1}(1-p)^{n+m-1-|\omega(a)|_1},
\qquad
\omega(a):=B_{m,n}^{-1}(a).
\]
Hence
$H_{\mu_{m,p}}(Y_0^{n-1})=(n+m-1)h_2(p)$, the excess entropy equals
$(m-1)h_2(p)$, and the coordinate process under $\mu_{m,p}$ is exact
$m$-step Markov.  The random local Fibonacci residue process
$R_t:=(\sum_{j=0}^{m-1}X_{t+j}F_{j+2})\bmod F_{m+2}$ under i.i.d.\
Bernoulli$(p)$ input is itself an exact $m$-step Markov process.
\item\label{item:main-inverse}
\textup{(Sliding block inverses.)}
Exact recovery from overlapping normalized windows is possible from any
length-$m$ observation window: for each $0\le r\le m-1$ there is a sliding
block inverse with memory $r$ and anticipation $m-1-r$.  At the one-sided
endpoint $r=m-1$, no smaller memory suffices.
\item\label{item:main-fischer}
\textup{(Fischer cover and presentation complexity.)}
Over the natural alphabet $X_m$, the canonical minimal right-resolving
presentation of $Y_m$ is the explicit suffix graph $G_m$, so it has exactly
$2^{m-1}$ states and $2^m$ edges.  Since $\lvert X_m\rvert=F_{m+2}$, the
natural deterministic presentation size grows exponentially relative to the
alphabet size.
\end{enumerate}
\end{theorem}

Taken together, these results show that a single integer controls the
finite-window arithmetic of the problem: it is the point at which local
Zeckendorf normalization becomes lossless on blocks, the sharp threshold for
local Fibonacci congruence rigidity, the exact Markov memory of the random
residue process, and the scale parameter of the natural deterministic
presentation; see again
Corollary~\ref{cor:local-complexity-scale-coincidence} and
Theorem~\ref{thm:natural-state-complexity-gap}.

These statements are proved in
Theorem~\ref{thm:block-bijection-threshold},
Corollary~\ref{cor:block-counts-threshold},
Corollary~\ref{cor:higher-block-identification},
Theorem~\ref{thm:Phi-m-conjugacy-threshold},
Theorem~\ref{thm:bernoulli-transport-cylinder-law},
Corollary~\ref{cor:exact-entropy-profile},
Theorem~\ref{thm:exact-markov-order-bernoulli},
Corollary~\ref{cor:Phi-m-memory-threshold-exact},
Theorem~\ref{thm:Y-m-exact-SFT-order},
Corollary~\ref{cor:uniform-synchronizing-length},
Proposition~\ref{prop:ambiguity-shell-nilpotency}, and
Corollary~\ref{cor:G-m-fischer-cover}, together with
Theorem~\ref{thm:natural-state-complexity-gap}.

\paragraph{Sketch of the reconstruction mechanism.}
The inverse is built in two steps.  First, the pairwise overlap maps
$\Theta_m$ compress adjacent folded symbols to the fold of their common raw
overlap; iterating this reduction through a block
$a_1\cdots a_m\in\mathcal L_m(Y_m)$ yields the central bit
$\Theta_{m,m}(a_1\cdots a_m)$.  Second, once this central bit and the folded
inner overlaps are known, the one-step left- and right-resolving properties of
$G_m$ force the neighboring bits uniquely, so an outward induction recovers the
entire raw lift of length $2m-1$.  The inverse family in Section~7 simply
records the same reconstruction after choosing which coordinate of the
recovered length-$m$ raw window is read as the output bit.

The paper is organized as follows.  Sections~2 and~3 develop finite-window
normalization and its rewrite description.  Section~4 proves compatibility in
the window length, Section~5 formulates the overlap-decoding problem,
Section~6 records the explicit finite-state presentation used later, and
Section~7 proves the block reconstruction theorem, derives the measure-theoretic
consequences, and then studies the natural right-resolving presentation of the
folded language.  Section~8 reformulates the main results in purely
number-theoretic terms: local Fibonacci congruence rigidity, representation
fingerprint separation, exact bit recovery from residue windows, and the
precise distribution and Markov order of random local residue vectors.



\section{Finite-window normalization}

We introduce the finite languages and the map $\Fold_m$.  The only
numeration-theoretic input needed later is Zeckendorf uniqueness, which
identifies the golden-mean language as the canonical target of the fold.

\subsection{Admissible words}

Throughout the paper the Fibonacci numbers are normalized by
\[
F_1=F_2=1,
\qquad
F_{k+2}=F_{k+1}+F_k
\quad (k\ge 1).
\]

For $m\ge 1$ and $s\in\{0,1\}$ let
\[
X_m^{(s)}:=\{w=w_1\cdots w_m\in X_m:w_m=s\}.
\]
Thus $X_m^{(0)}$ and $X_m^{(1)}$ split admissible words by their last symbol.

\begin{proposition}[Terminal recursion for admissible words]
\label{prop:stable-syntax-terminal-recursion}
For every $m\ge 1$, if
\[
v_m:=
\begin{pmatrix}
\lvert X_m^{(1)}\rvert\\[2pt]
\lvert X_m^{(0)}\rvert
\end{pmatrix},
\]
then
\[
v_{m+1}=N_\tau v_m,
\qquad
N_\tau=
\begin{pmatrix}
0&1\\[2pt]
1&1
\end{pmatrix}.
\]
\end{proposition}

\begin{proof}
If $w\in X_{m+1}$ ends in $1$, then its penultimate symbol must be $0$.
Deleting the last symbol therefore gives a bijection
\[
X_{m+1}^{(1)}\longleftrightarrow X_m^{(0)},
\]
so $\lvert X_{m+1}^{(1)}\rvert=\lvert X_m^{(0)}\rvert$.

If $w\in X_{m+1}$ ends in $0$, then deleting the last symbol gives an
arbitrary word in $X_m$, and appending $0$ to a word in $X_m$ never creates
the forbidden pattern $11$.  Hence
\[
\lvert X_{m+1}^{(0)}\rvert
=
\lvert X_m\rvert
=
\lvert X_m^{(0)}\rvert+\lvert X_m^{(1)}\rvert.
\]
Combining the two identities gives the displayed matrix recursion.
\end{proof}

\begin{corollary}[Fibonacci counting]
\label{cor:stable-syntax-counting}
For every $m\ge 1$ one has
\[
\lvert X_m\rvert=F_{m+2}.
\]
\end{corollary}

\begin{proof}
The recursion in Proposition~\ref{prop:stable-syntax-terminal-recursion} is
the Fibonacci recursion with the initial vector
\[
v_1=
\begin{pmatrix}
1\\[2pt]
1
\end{pmatrix}.
\]
Therefore
\[
\lvert X_m\rvert
=
\lvert X_m^{(0)}\rvert+\lvert X_m^{(1)}\rvert
=
F_{m+2}.
\qedhere
\]
\end{proof}

\subsection{Fibonacci value and the fold}

Let
\[
\Xfin
:=
\left\{
c=(c_1,c_2,\ldots)\in\{0,1\}^{\NN}:
\supp(c)\text{ finite and }c_kc_{k+1}=0\text{ for all }k
\right\}.
\]
For $c\in\Xfin$ define
\[
\Val(c):=\sum_{k\ge 1} c_kF_{k+1}.
\]
Zeckendorf's theorem gives, for every $N\in\NN$, a unique $Z(N)\in \Xfin$
satisfying $\Val(Z(N))=N$ \cite{Zeckendorf1972}.

For a finite binary window $\omega=\omega_1\cdots\omega_m\in\Omega_m$ define
its Fibonacci value by
\[
N(\omega):=\sum_{k=1}^m \omega_kF_{k+1}.
\]
Let
\[
\pi_m:\Xfin\to X_m,
\qquad
\pi_m(c_1,c_2,\ldots):=(c_1,\ldots,c_m),
\]
be the length-$m$ prefix projection.

\begin{definition}[Finite-window fold]
\label{def:fold-word}
For $m\ge 1$ define
\[
\Fold_m:\Omega_m\to X_m,
\qquad
\Fold_m(\omega):=\pi_m(Z(N(\omega))).
\]
\end{definition}

\begin{definition}[Visible value]
\label{def:visible-value}
For $x=x_1\cdots x_m\in X_m$ define
\[
V_m(x):=\sum_{k=1}^m x_kF_{k+1}.
\]
\end{definition}

\begin{proposition}[Visible-value parametrization of $X_m$]
\label{prop:visible-value-parametrization}
For every $m\ge 1$, the map
\[
V_m:X_m\to \{0,1,\ldots,F_{m+2}-1\}
\]
is a bijection.
\end{proposition}

\begin{proof}
If $x\in X_m$, then
\[
(x_1,\ldots,x_m,0,0,\ldots)\in \Xfin
\]
is already a Zeckendorf digit string, so by uniqueness it is the Zeckendorf
representation of the integer $V_m(x)$.  In particular $V_m(x)<F_{m+2}$,
because otherwise the Zeckendorf representation would have to use the weight
$F_{m+2}$ or a larger one.

Conversely, let $0\le r<F_{m+2}$.  The Zeckendorf representation $Z(r)$ cannot
use the weight $F_{m+2}$ or any larger weight, so it has the form
\[
Z(r)=(x_1,\ldots,x_m,0,0,\ldots)
\]
for a unique $x\in X_m$.  Then $V_m(x)=r$.  Thus $V_m$ is bijective.
\end{proof}

\begin{corollary}[Unique overflow decomposition]
\label{cor:visible-value-overflow}
For every $m\ge 1$ and every $\omega\in\Omega_m$, there exists a unique
$b(\omega)\in\{0,1\}$ such that
\[
N(\omega)=V_m(\Fold_m(\omega))+b(\omega)F_{m+2}.
\]
Equivalently, $\Fold_m(\omega)$ is the unique admissible word in $X_m$ whose
visible value is congruent to $N(\omega)$ modulo $F_{m+2}$.
\end{corollary}

\begin{proof}
Since
\[
0\le N(\omega)\le \sum_{k=1}^m F_{k+1}=F_{m+3}-2<2F_{m+2},
\]
there is a unique $b(\omega)\in\{0,1\}$ such that
\[
r:=N(\omega)-b(\omega)F_{m+2}
\]
lies in $\{0,1,\ldots,F_{m+2}-1\}$.  By
Proposition~\ref{prop:visible-value-parametrization}, there is a unique
$x\in X_m$ with $V_m(x)=r$.  The first $m$ digits of $Z(N(\omega))$ are
precisely this $x$, so $x=\Fold_m(\omega)$.
\end{proof}

For convenience, Table~\ref{tab:value-maps} summarizes the value maps used
throughout the paper.

\begin{table}[t]
\centering
\small
\begin{tabular}{lll}
\toprule
symbol & domain & role \\
\midrule
$N(\omega)$ & $\Omega_m$ & full Fibonacci value of the raw window $\omega$ \\
$V_m(x)$ & $X_m$ & visible value carried by the admissible word $x$ \\
$b(\omega)$ & $\Omega_m$ & overflow bit beyond the visible range $[0,F_{m+2})$ \\
$\Val(a)$ & $\mathcal D$ & value of a finitely supported digit configuration \\
\bottomrule
\end{tabular}
\caption{Value maps used throughout the paper.  The visible value stays within
the length-$m$ window, while the overflow bit records whether one copy of
$F_{m+2}$ has escaped beyond it.}
\label{tab:value-maps}
\end{table}



\section{A local rewrite description}

In this section we describe $\Fold_m$ by local rewriting on finitely supported
configurations.  This gives a concrete normalization procedure before we pass
to inverse limits and sequence-level dynamics.

\subsection{Value-preserving local rules}

Let
\[
\mathcal D
:=
\left\{
a=(a_1,a_2,\ldots):
a_k\in\NN,\ \supp(a)\text{ finite}
\right\},
\]
and define
\[
\Val(a):=\sum_{k\ge 1} a_kF_{k+1}.
\]
We write $e_k$ for the $k$th standard basis vector and
\[
\iota_m(\omega):=(\omega_1,\ldots,\omega_m,0,0,\ldots)\in\mathcal D
\]
for the natural embedding of a window into $\mathcal D$.

The rewrite system acts on $\mathcal D$ through the following local rules:
\begin{align}
e_k+e_{k+1}&\longrightarrow e_{k+2} && (k\ge 1), \label{eq:rewrite-adjacent}\\
2e_1&\longrightarrow e_2, \label{eq:rewrite-base1}\\
2e_2&\longrightarrow e_1+e_3, \label{eq:rewrite-base2}\\
2e_k&\longrightarrow e_{k-2}+e_{k+1} && (k\ge 3). \label{eq:rewrite-duplicate}
\end{align}
Each rule is local and preserves $\Val$ because it is exactly the Fibonacci
identity
\[
F_{k+1}+F_{k+2}=F_{k+3},
\qquad
2F_2=F_3,
\qquad
2F_3=F_2+F_4,
\qquad
2F_{k+1}=F_{k-1}+F_{k+2}.
\]

\begin{lemma}[Fibonacci rigidity of value-preserving local weights]
\label{lem:fold-local-weight-rigidity-fibonacci}
Suppose a positive weight sequence $(w_k)_{k\ge 1}$ has the property that each
rewrite rule \eqref{eq:rewrite-adjacent}--\eqref{eq:rewrite-duplicate}
preserves the weighted sum
\[
\Val_w(a):=\sum_{k\ge 1} a_kw_k
\qquad (a\in\mathcal D).
\]
Then
\[
w_{k+2}=w_{k+1}+w_k
\qquad (k\ge 1).
\]
In particular, under the normalization $w_1=F_2=1$ and $w_2=F_3=2$, one has
$w_k=F_{k+1}$ for all $k\ge 1$.
\end{lemma}

\begin{proof}
Preservation of \eqref{eq:rewrite-adjacent} gives
\[
w_k+w_{k+1}=w_{k+2}
\qquad (k\ge 1),
\]
which is exactly the Fibonacci recursion.  The remaining rules are then
automatic consequences of that recursion:
\[
2w_1=w_2,\qquad 2w_2=w_1+w_3,\qquad 2w_k=w_{k-2}+w_{k+1}\ \ (k\ge 3).
\]
With the stated normalization, the recursion has the unique solution
$w_k=F_{k+1}$.
\end{proof}

\begin{lemma}[Local weight rigidity]
\label{lem:local-weight-rigidity}
Let
\[
A(a):=\sum_{k\ge 1} a_k,
\qquad
B(a):=\sum_{k\ge 1} ka_k,
\qquad
C(a):=a_2,
\]
and order triples lexicographically.  Then every rewrite step
$a\to a'$ generated by \eqref{eq:rewrite-adjacent}--\eqref{eq:rewrite-duplicate}
strictly decreases
\[
\mu(a):=(A(a),B(a),C(a)).
\]
In particular the rewrite system is strongly terminating.
\end{lemma}

\begin{proof}
Rules \eqref{eq:rewrite-adjacent} and \eqref{eq:rewrite-base1} reduce the total
digit count $A$ by one.  Rule \eqref{eq:rewrite-duplicate} for $k\ge 3$
preserves $A$ but changes $B$ by
\[
2k-\bigl((k-2)+(k+1)\bigr)=1,
\]
so $B$ decreases by one.  Rule \eqref{eq:rewrite-base2} preserves both $A$ and
$B$ but decreases $C=a_2$ by two.  Hence every rewrite step strictly decreases
$\mu(a)$.
\end{proof}

\subsection{Canonical normal forms}

\begin{proposition}[Termination, confluence, and normal forms]
\label{prop:rewrite-newman}
The rewrite system on $\mathcal D$ generated by
\eqref{eq:rewrite-adjacent}--\eqref{eq:rewrite-duplicate} is strongly
terminating and confluent.  Every $a\in\mathcal D$ therefore has a unique
normal form
\[
\NF(a)\in\Xfin
\]
satisfying
\[
\Val(\NF(a))=\Val(a).
\]
\end{proposition}

\begin{proof}
Strong termination is Lemma~\ref{lem:local-weight-rigidity}.  A configuration
is irreducible exactly when every digit is in $\{0,1\}$ and no adjacent pair
$11$ occurs; this is precisely the definition of $\Xfin$.

Fix $a\in\mathcal D$.  By termination, every maximal reduction sequence from
$a$ ends in some irreducible configuration $x\in\Xfin$.  Since each rewrite
preserves $\Val$, every such terminal configuration satisfies
\[
\Val(x)=\Val(a).
\]
If $x$ and $y$ are two terminal reductions of $a$, then $x,y\in\Xfin$ are
irreducible configurations with the same Fibonacci value.  Zeckendorf
uniqueness therefore forces $x=y$.  Thus every $a$ has a unique irreducible
descendant.  For a terminating rewrite system, uniqueness of irreducible
descendants is equivalent to confluence, so the system is confluent and the
normal form $\NF(a)$ is well defined.
\end{proof}

\begin{corollary}[Order-independent implementation of the fold]
\label{cor:fold-order-independent}
For every $\omega\in\Omega_m$,
\[
\Fold_m(\omega)=\pi_m(\NF(\iota_m(\omega))),
\]
and the right-hand side is independent of the order in which rewrites are
applied.
\end{corollary}

\begin{proof}
By Proposition~\ref{prop:rewrite-newman}, the normal form of $\iota_m(\omega)$
is the unique admissible Fibonacci digit string with value $N(\omega)$, namely
$Z(N(\omega))$.  Taking the first $m$ digits gives the claim.
\end{proof}

\begin{remark}[Two illustrative normalizations]
\label{rem:two-fold-examples}
The following two examples illustrate the order-independence of the rewriting
and the appearance of the overflow bit.

For $\omega=1110\in\Omega_4$ one may reduce
\[
\iota_4(\omega)=e_1+e_2+e_3\to 2e_3\to e_1+e_4,
\]
so $\NF(\iota_4(\omega))=\iota_4(1001)$ and
\[
\Fold_4(1110)=1001.
\]
If instead $\omega=1101\in\Omega_4$, then
\[
\iota_4(\omega)=e_1+e_2+e_4\to e_3+e_4\to e_5,
\]
hence
\[
\Fold_4(1101)=0000,
\qquad
N(1101)=F_6=F_{4+2}.
\]
This is the smallest example in which the visible word is zero while the whole
value survives as the overflow bit from
Corollary~\ref{cor:visible-value-overflow}.
\end{remark}

\begin{proposition}[Basic properties of $\Fold_m$]
\label{prop:fold-basic}
For every $m\ge 1$, the map $\Fold_m:\Omega_m\to X_m$ is well defined and
satisfies:
\begin{enumerate}[label=\textup{(\roman*)}]
\item $\Fold_m(\omega)\in X_m$ for every $\omega\in\Omega_m$;
\item $\Fold_m(x)=x$ for every $x\in X_m$;
\item $\Fold_m:\Omega_m\twoheadrightarrow X_m$ is surjective.
\end{enumerate}
\end{proposition}

\begin{proof}
Property (i) is immediate from Definition~\ref{def:fold-word}.  For (ii), if
$x\in X_m$ then $x$ is already an admissible Zeckendorf prefix, so
$Z(N(x))=(x_1,\ldots,x_m,0,0,\ldots)$ and hence $\Fold_m(x)=x$.  Surjectivity
follows from (ii): every $x\in X_m$ is a fixed point of $\Fold_m$.
\end{proof}



\section{Compatibility in the window length}

We next compare the folds at different window lengths.  The results of this
section show that admissible restriction and finite-window normalization are
compatible with one another.

\subsection{Restriction maps}

\begin{definition}[Prefix restrictions]
\label{def:prefix-restriction}
For integers $n\ge m\ge 1$ define
\[
\pi_{n\to m}:X_n\to X_m,
\qquad
\pi_{n\to m}(x_1\cdots x_n):=x_1\cdots x_m.
\]
\end{definition}

\begin{lemma}[Functoriality of admissible restrictions]
\label{lem:pi-functorial}
If $n\ge \ell\ge m\ge 1$, then
\[
\pi_{n\to m}=\pi_{\ell\to m}\circ \pi_{n\to \ell}.
\]
Moreover, for every $c\in \Xfin$ and every $n\ge m$,
\[
\pi_{n\to m}(\pi_n(c))=\pi_m(c).
\]
\end{lemma}

\begin{proof}
All maps involved are literal prefix truncations, so both identities are
immediate.
\end{proof}

\begin{theorem}[Inverse-limit identification]
\label{thm:inverse-limit-golden}
Let
\[
\Xinf
:=
\{c=(c_1,c_2,\ldots)\in\{0,1\}^{\NN}:c_kc_{k+1}=0\text{ for all }k\ge 1\}.
\]
Then there is a natural bijection
\[
\Xinf\cong \varprojlim (X_m,\pi_{m+1\to m}).
\]
\end{theorem}

\begin{proof}
Given $c\in \Xinf$, its prefix family
\[
(\pi_m(c))_{m\ge 1}
\]
is compatible under the restriction maps, hence defines an element of the
inverse limit.  Conversely, given a compatible family
\[
(x^{(m)})_{m\ge 1}\in \varprojlim (X_m,\pi_{m+1\to m}),
\]
define $c_k:=x_k^{(m)}$ for any $m\ge k$.  Compatibility ensures that this is
independent of the choice of $m$.  Since every finite prefix $x^{(m)}$ avoids
$11$, the resulting sequence $c$ lies in $\Xinf$.  The two constructions are
inverse to one another.
\end{proof}

\begin{remark}
The object in Theorem~\ref{thm:inverse-limit-golden} is one sided because the
resolution tower is indexed by prefix length.  It should not be confused with
the two-sided image shift $Y_m$ introduced later at the sequence level.
\end{remark}

\subsection{Compatibility with raw windows}

If one truncates a raw window directly, folding does not in general commute
with restriction.  Let
\[
r_{m+1\to m}:\Omega_{m+1}\to\Omega_m,
\qquad
r_{m+1\to m}(\omega_1\cdots\omega_{m+1})=\omega_1\cdots\omega_m.
\]
For $m=1$ and $\omega=11$ one has
\[
(\Fold_1\circ r_{2\to 1})(11)=\Fold_1(1)=1,
\]
whereas
\[
(\pi_{2\to 1}\circ \Fold_2)(11)=0,
\]
since $N(11)=F_2+F_3=F_4$ and hence $\Fold_2(11)=00$.

The correct raw-level restriction is therefore the folded one.

\begin{definition}[Fold-compatible raw restriction]
\label{def:fold-compatible-restriction}
For $m\ge 1$ define
\[
\widetilde r_{m+1\to m}:\Omega_{m+1}\to\Omega_m,
\qquad
\widetilde r_{m+1\to m}(\omega):=\pi_{m+1\to m}(\Fold_{m+1}(\omega)).
\]
\end{definition}

\begin{proposition}[Commuting diagram for fold-compatible restriction]
\label{prop:fold-omega-commute}
For every $m\ge 1$ and every $\omega\in\Omega_{m+1}$,
\[
\Fold_m(\widetilde r_{m+1\to m}(\omega))
=
\pi_{m+1\to m}(\Fold_{m+1}(\omega)).
\]
Equivalently, the diagram
\[
\begin{CD}
\Omega_{m+1} @>{\Fold_{m+1}}>> X_{m+1}\\
@V{\widetilde r_{m+1\to m}}VV @VV{\pi_{m+1\to m}}V\\
\Omega_m @>{\Fold_m}>> X_m
\end{CD}
\]
commutes.
\end{proposition}

\begin{proof}
By definition,
\[
\widetilde r_{m+1\to m}(\omega)=\pi_{m+1\to m}(\Fold_{m+1}(\omega))\in X_m.
\]
Applying Proposition~\ref{prop:fold-basic}(ii) to this admissible word gives
\[
\Fold_m(\widetilde r_{m+1\to m}(\omega))
=
\widetilde r_{m+1\to m}(\omega)
=
\pi_{m+1\to m}(\Fold_{m+1}(\omega)).
\qedhere
\]
\end{proof}

\begin{remark}[Two raw-level restriction examples]
\label{rem:restriction-examples}
The discrepancy between direct restriction and folded restriction already
appears at the smallest scale, where $m=1$ and $\omega=11$:
\[
(\Fold_1\circ r_{2\to 1})(11)=1,
\qquad
(\Fold_1\circ \widetilde r_{2\to 1})(11)=0.
\]
At the next nontrivial scale one may take $m=3$ and $\omega=1101$.  Then
\[
r_{4\to 3}(1101)=110,
\qquad
\Fold_3(110)=001,
\]
whereas Remark~\ref{rem:two-fold-examples} gives $\Fold_4(1101)=0000$, so
\[
\widetilde r_{4\to 3}(1101)=000,
\qquad
\Fold_3(\widetilde r_{4\to 3}(1101))=000.
\]
Thus the naive truncation retains a visible carry that disappears once one
first folds at the larger resolution and only then restricts.
\end{remark}



\section{Overlapping normalized windows}

We now pass to bi-infinite sequences by sliding the window along the integers.
This places the construction in the standard framework of factor maps between
shift spaces.

Let $\sigma$ denote the left shift on both $\{0,1\}^{\ZZ}$ and $X_m^{\ZZ}$.

\begin{definition}[Sequence-level fold]
\label{def:Phi-m}
For $m\ge 1$ define
\[
\Phi_m:\{0,1\}^{\ZZ}\to X_m^{\ZZ},
\qquad
(\Phi_m(s))_t:=\Fold_m(s_t s_{t+1}\cdots s_{t+m-1}).
\]
\end{definition}

\begin{proposition}[Finite radius]
\label{prop:Phi-m-radius}
The map $\Phi_m$ is a sliding block code with memory $0$ and anticipation
$m-1$.
\end{proposition}

\begin{proof}
By definition, the symbol $(\Phi_m(s))_t$ depends only on the coordinates
$s_t,\ldots,s_{t+m-1}$.
\end{proof}

\begin{proposition}[Shift equivariance]
\label{prop:Phi-m-equivariant}
For every $m\ge 1$,
\[
\Phi_m\circ \sigma=\sigma\circ \Phi_m.
\]
\end{proposition}

\begin{proof}
Both sides evaluate $\Fold_m$ on the same translated length-$m$ window at each
time $t$.
\end{proof}

\begin{proposition}[Factor property on subshifts]
\label{prop:Phi-m-factor}
Let $S\subset \{0,1\}^{\ZZ}$ be a subshift.  Then $\Phi_m|_S:S\to X_m^{\ZZ}$ is
continuous, shift commuting, and onto its image
\[
Y_m(S):=\Phi_m(S).
\]
Moreover $Y_m(S)$ is a subshift of $X_m^{\ZZ}$, so
\[
\Phi_m|_S:S\twoheadrightarrow Y_m(S)
\]
is a factor map.
\end{proposition}

\begin{proof}
Continuity follows from Proposition~\ref{prop:Phi-m-radius}, and
shift-equivariance from Proposition~\ref{prop:Phi-m-equivariant}.  Since $S$ is
compact and shift invariant, its image under the continuous equivariant map
$\Phi_m$ is again compact and shift invariant.  By definition, the restriction
is onto that image.
\end{proof}

\begin{definition}[The image shift]
\label{def:Y-m}
We write
\[
Y_m:=\Phi_m(\{0,1\}^{\ZZ})\subset X_m^{\ZZ}.
\]
\end{definition}



\section{Finite-state structure of normalized windows}

This section records the explicit suffix graph attached to the folded language.
It will be used throughout the reconstruction arguments in the final section.

\subsection{The graph \texorpdfstring{$G_m$}{Gm}}

\begin{definition}[Suffix-state graph]
\label{def:G-m}
For $m\ge 1$, let
\[
V_m:=\{0,1\}^{m-1}.
\]
For each $v=v_1\cdots v_{m-1}\in V_m$ and each $b\in\{0,1\}$ define an edge
\[
e=(v,b):v\to v':=v_2\cdots v_{m-1}b,
\]
labeled by
\[
\ell_m(e):=\Fold_m(v_1\cdots v_{m-1}b)\in X_m.
\]
The resulting finite labeled graph is denoted
\[
G_m=(V_m,E_m,\ell_m).
\]
It has $\lvert V_m\rvert=2^{m-1}$ vertices and $\lvert E_m\rvert=2^m$ edges.
\end{definition}

\begin{theorem}[Suffix-graph presentation of the folded language]
\label{thm:Phi-m-sofic-graph}
Let $Z_m\subset E_m^{\ZZ}$ be the two-sided edge shift of $G_m$, and let
\[
\mathcal L_m:Z_m\to X_m^{\ZZ}
\]
be the edge-label map.  Then
\[
Y_m=\mathcal L_m(Z_m).
\]
In particular $G_m$ presents the folded language $Y_m$.
\end{theorem}

\begin{proof}
Take any $s\in\{0,1\}^{\ZZ}$ and define
\[
v_t:=s_t\cdots s_{t+m-2}\in V_m,
\qquad
e_t:=(v_t,s_{t+m-1})\in E_m.
\]
The compatibility of successive windows implies that $(e_t)_{t\in\ZZ}$ is a
bi-infinite path in $G_m$, and by construction
\[
(\mathcal L_m(e))_t
=
\ell_m(e_t)
=
\Fold_m(s_t\cdots s_{t+m-1})
=
(\Phi_m(s))_t.
\]
Hence $Y_m\subset \mathcal L_m(Z_m)$.

Conversely, let $e=(e_t)_{t\in\ZZ}\in Z_m$.  Writing each edge as
$e_t=(v_t,b_t)$, path compatibility says that $v_{t+1}$ is obtained from $v_t$
by deleting its first symbol and appending $b_t$.  Therefore the family of
suffixes $v_t$ and appended bits $b_t$ determines a unique binary sequence
$s\in\{0,1\}^{\ZZ}$ such that
\[
s_t\cdots s_{t+m-1}=v_tb_t
\]
for all $t$.  Then
\[
(\Phi_m(s))_t
=
\Fold_m(v_tb_t)
=
\ell_m(e_t)
=
(\mathcal L_m(e))_t.
\]
Thus $\mathcal L_m(Z_m)\subset Y_m$, proving equality.
\end{proof}

\begin{corollary}[The image shift is sofic]
\label{prop:Y-m-sofic}
For every $m\ge 1$, the image shift $Y_m$ is a sofic shift.
\end{corollary}

\begin{proof}
This is immediate from Theorem~\ref{thm:Phi-m-sofic-graph}, which gives an
explicit finite labeled graph presenting $Y_m$.
\end{proof}

\subsection{Deterministic structure of the presentation}

We begin with the one-step resolving properties of the graph $G_m$.

\begin{proposition}[One-step distinguishability]
\label{prop:one-step-right-resolving}
For every $m\ge 2$ and every $v\in\{0,1\}^{m-1}$,
\[
\Fold_m(v0)\neq \Fold_m(v1).
\]
Consequently $G_m$ is right-resolving: the two outgoing edges from any fixed
vertex carry distinct labels.
\end{proposition}

\begin{proof}
Let $\omega_0:=v0$ and $\omega_1:=v1$.  Then
\[
N(\omega_1)-N(\omega_0)=F_{m+1}.
\]
If $\Fold_m(\omega_0)=\Fold_m(\omega_1)$, then
Corollary~\ref{cor:visible-value-overflow} implies
\[
\begin{aligned}
N(\omega_1)-N(\omega_0)
&=
\bigl(b(\omega_1)-b(\omega_0)\bigr)F_{m+2}\\
&\in\{-F_{m+2},0,F_{m+2}\},
\end{aligned}
\]
because the visible value $V_m(\Fold_m(\omega_i))$ would be the same for
$i=0,1$ and only the overflow bit could differ.  This contradicts
$F_{m+1}>0$ and $F_{m+1}\neq F_{m+2}$.  Hence the two labels are distinct.
\end{proof}

\begin{proposition}[One-step left-resolving property]
\label{prop:one-step-left-resolving}
For every $m\ge 2$ and every $v\in\{0,1\}^{m-1}$,
\[
\Fold_m(0v)\neq \Fold_m(1v).
\]
Consequently the two incoming edges to any fixed vertex of $G_m$ carry
distinct labels.
\end{proposition}

\begin{proof}
Let $\omega_0:=0v$ and $\omega_1:=1v$.  Then
\[
N(\omega_1)-N(\omega_0)=F_2=1.
\]
If $\Fold_m(\omega_0)=\Fold_m(\omega_1)$, then
Corollary~\ref{cor:visible-value-overflow} implies
\[
N(\omega_1)-N(\omega_0)
=
\bigl(b(\omega_1)-b(\omega_0)\bigr)F_{m+2}
\in\{-F_{m+2},0,F_{m+2}\},
\]
which is impossible because $F_{m+2}\ge 2$ for $m\ge 2$.  Hence the labels are
distinct.  Applied to the two incoming edges of a fixed vertex
$v=v_1\cdots v_{m-1}$, this says the edges labeled by the windows
$0v_1\cdots v_{m-1}$ and $1v_1\cdots v_{m-1}$ are distinguished.
\end{proof}

\begin{corollary}[Entropy of the image shift]
\label{cor:Phi-m-entropy}
For every $m\ge 2$,
\[
h_{\mathrm{top}}(Y_m)=\log 2.
\]
\end{corollary}

\begin{proof}
Since $\Phi_m$ is a factor of the full two-shift,
\[
h_{\mathrm{top}}(Y_m)\le \log 2.
\]
On the other hand, every vertex of $G_m$ has exactly two outgoing edges, so the
number of length-$n$ paths starting from any fixed vertex is $2^n$.  By
Proposition~\ref{prop:one-step-right-resolving}, two such paths cannot have the
same label sequence.  Hence $Y_m$ has at least $2^n$ allowed blocks of length
$n$, so $h_{\mathrm{top}}(Y_m)\ge \log 2$.
\end{proof}

\subsection{Fibres and synchronizing blocks}

For $m\ge 1$ and $0\le r\le \lfloor m/2\rfloor$ define
\[
z_{m,0}:=0^m,
\qquad
z_{m,r}:=0^{m-2r}11(01)^{r-1}
\quad (1\le r\le \lfloor m/2\rfloor).
\]

\begin{proposition}[The zero fibre of $0^m$]
\label{prop:zero-fiber-explicit}
For every $m\ge 1$,
\[
\Fold_m^{-1}(0^m)=\{z_{m,r}:0\le r\le \lfloor m/2\rfloor\}.
\]
In particular,
\[
\bigl|\Fold_m^{-1}(0^m)\bigr|=\lfloor m/2\rfloor+1.
\]
\end{proposition}

\begin{proof}
By Corollary~\ref{cor:visible-value-overflow},
\[
\Fold_m(\omega)=0^m
\quad\Longleftrightarrow\quad
N(\omega)\in\{0,F_{m+2}\}.
\]
The case $N(\omega)=0$ forces $\omega=0^m$, so it remains to classify
\[
A_m:=\{\omega\in\Omega_m:N(\omega)=F_{m+2}\}.
\]

For $m=1$ the set $A_1$ is empty, since the only available values are $0$ and
$F_2=1<F_3$.  Fix now $m\ge 2$ and let $\omega\in A_m$.  Then
\[
\sum_{k=1}^{m-1}\omega_kF_{k+1}\le \sum_{k=1}^{m-1}F_{k+1}=F_{m+2}-2<F_{m+2},
\]
so necessarily $\omega_m=1$.  If also $\omega_{m-1}=1$, then
\[
\sum_{k=1}^{m-2}\omega_kF_{k+1}
=
F_{m+2}-F_{m+1}-F_m
=
0,
\]
hence $\omega=0^{m-2}11$.  If instead $\omega_{m-1}=0$, write $\omega=u01$
with $u\in\Omega_{m-2}$.  Then
\[
N(u)=F_{m+2}-F_{m+1}=F_m,
\]
so $u\in A_{m-2}$.  Therefore, for every $m\ge 2$,
\[
A_m=\{0^{m-2}11\}\cup\{u01:u\in A_{m-2}\}.
\]

Starting from $A_1=\varnothing$ and $A_2=\{11\}$ and unwinding this recursion
gives
\[
A_m=\{z_{m,r}:1\le r\le \lfloor m/2\rfloor\}.
\]
Adding the zero word yields the stated formula for $\Fold_m^{-1}(0^m)$.  The
cardinality follows immediately.
\end{proof}

\begin{definition}[Alternating extremal words]
\label{def:alternating-extremal-words}
For each $m\ge 1$, let $a_m,b_m\in X_m$ denote the unique alternating
admissible words of length $m$ ending respectively in $0$ and $1$.  Thus
\[
a_1=0,
\qquad
b_1=1,
\]
and for $m\ge 2$,
\[
a_m=b_{m-1}0,
\qquad
b_m=a_{m-1}1.
\]
\end{definition}

\begin{lemma}[Extremal values of the alternating words]
\label{lem:alternating-extremal-values}
For every $m\ge 1$ one has
\[
V_m(a_m)=F_{m+1}-1,
\qquad
V_m(b_m)=F_{m+2}-1.
\]
\end{lemma}

\begin{proof}
We argue by induction on $m$.  For $m=1$ the claim is
\[
V_1(a_1)=0=F_2-1,
\qquad
V_1(b_1)=1=F_3-1.
\]

Assume now that $m\ge 2$ and the claim is known for $m-1$.  Since $a_m=b_{m-1}0$,
one has
\[
V_m(a_m)=V_{m-1}(b_{m-1})=F_{m+1}-1.
\]
Likewise, using $b_m=a_{m-1}1$,
\[
V_m(b_m)=V_{m-1}(a_{m-1})+F_{m+1}
=(F_m-1)+F_{m+1}
=F_{m+2}-1.
\]
\end{proof}

\begin{proposition}[Singleton fibres of the alternating symbols]
\label{prop:alternating-singleton-fibres}
For every $m\ge 1$,
\[
\Fold_m^{-1}(a_m)=\{a_m\},
\qquad
\Fold_m^{-1}(b_m)=\{b_m\}.
\]
\end{proposition}

\begin{proof}
The case $m=1$ is immediate.  Fix $m\ge 2$.

Suppose first that $\Fold_m(\omega)=a_m$.  By
Lemma~\ref{lem:alternating-extremal-values} and
Corollary~\ref{cor:visible-value-overflow},
\[
N(\omega)\in\{F_{m+1}-1,\;F_{m+1}+F_{m+2}-1\}.
\]
But
\[
F_{m+1}+F_{m+2}-1=F_{m+3}-1>F_{m+3}-2\ge N(\omega),
\]
so necessarily $N(\omega)=F_{m+1}-1$.  In particular $N(\omega)<F_{m+1}$, and
therefore $\omega_m=0$.  Writing $\omega=u0$ with $u\in\Omega_{m-1}$, we get
\[
N(u)=F_{m+1}-1.
\]
By the induction hypothesis, $u=b_{m-1}$.  Hence $\omega=b_{m-1}0=a_m$.

Now suppose that $\Fold_m(\omega)=b_m$.  Again by
Lemma~\ref{lem:alternating-extremal-values} and
Corollary~\ref{cor:visible-value-overflow},
\[
N(\omega)\in\{F_{m+2}-1,\;2F_{m+2}-1\}.
\]
Since
\[
2F_{m+2}-1
>
F_{m+2}+F_{m+1}-2
=
F_{m+3}-2
\ge
N(\omega),
\]
one must have $N(\omega)=F_{m+2}-1$.  The contribution of the first $m-1$
digits is at most
\[
\sum_{k=1}^{m-1}F_{k+1}=F_{m+2}-2,
\]
so necessarily $\omega_m=1$.  Writing $\omega=u1$ with $u\in\Omega_{m-1}$, we
obtain
\[
N(u)=F_{m+2}-1-F_{m+1}=F_m-1.
\]
By the induction hypothesis, $u=a_{m-1}$.  Therefore
\[
\omega=a_{m-1}1=b_m.
\]
\end{proof}

\begin{theorem}[Single-symbol synchronization]
\label{thm:alternating-single-symbol-synchronizing}
For every $m\ge 1$, each of the symbols $a_m$ and $b_m$ labels exactly one edge
of $G_m$.  Consequently both $a_m$ and $b_m$ are synchronizing words for the
presentation $G_m$ of $Y_m$.
\end{theorem}

\begin{proof}
An edge of $G_m$ is labeled by $x\in X_m$ precisely when its defining raw word
lies in $\Fold_m^{-1}(x)$.  Proposition~\ref{prop:alternating-singleton-fibres}
shows that $\Fold_m^{-1}(a_m)$ and $\Fold_m^{-1}(b_m)$ are singletons.  Hence
each of $a_m$ and $b_m$ occurs on exactly one edge of $G_m$.  Any label carried
by exactly one edge is synchronizing, since reading that label collapses the
candidate vertex set to the terminal vertex of the unique edge.
\end{proof}

\begin{corollary}[The image language has synchronizing symbols]
\label{cor:Y-m-synchronized}
For every $m\ge 1$, the symbols $a_m$ and $b_m$ are synchronizing for the
presentation $G_m$.  In particular, the image language $Y_m$ is synchronized.
\end{corollary}

\begin{proof}
By Proposition~\ref{prop:one-step-right-resolving}, the graph $G_m$ is a
right-resolving presentation of $Y_m$.  By
Theorem~\ref{thm:alternating-single-symbol-synchronizing}, that presentation
admits a synchronizing symbol.  Hence $Y_m$ is synchronized; see
\cite[Ch.~3]{LindMarcus1995} or \cite[Ch.~4]{Kitchens1998SymbolicDynamics}.
\end{proof}

\subsection{A lower bound for one-sided inversion}

\begin{proposition}[A lower bound for one-sided inversion]
\label{prop:one-sided-inverse-memory-lower-bound}
Let $m\ge 3$, and define finite binary words
\[
u_m:=0^m10^{m-2},
\qquad
v_m:=0^{m-2}110^{m-1}
\]
of length $2m-1$.  Let $U_m,V_m\in\{0,1\}^{\ZZ}$ be their zero extensions on
the coordinates $\{1,\ldots,2m-1\}$.  Also define
\[
e_{m,0}:=0^m,
\qquad
e_{m,j}:=0^{m-j}10^{j-1}
\quad (1\le j\le m-1).
\]
Then
\[
(\Phi_m(U_m))_t=(\Phi_m(V_m))_t=e_{m,t-1}
\qquad (1\le t\le m-1),
\]
while
\[
(\Phi_m(U_m))_m=010^{m-2}\neq 10^{m-1}=(\Phi_m(V_m))_m.
\]
Consequently, if
\[
\Psi:Y_m\to\{0,1\}^{\ZZ}
\]
is a sliding block code with memory $q$ and anticipation $0$ such that
\[
\Psi\circ\Phi_m=\mathrm{id}_{\{0,1\}^{\ZZ}},
\]
then necessarily $q\ge m-1$.
\end{proposition}

\begin{proof}
For $U_m$, the $t$th raw window is $e_{m,t-1}\in X_m$ for every
$1\le t\le m-1$, so
\[
(\Phi_m(U_m))_t=e_{m,t-1}
\qquad (1\le t\le m-1).
\]

For $V_m$, the first raw window is $0^{m-2}11=z_{m,1}$, so
Proposition~\ref{prop:zero-fiber-explicit} gives
\[
(\Phi_m(V_m))_1=0^m=e_{m,0}.
\]
For $2\le t\le m-1$, the $t$th raw window of $V_m$ is
\[
0^{m-t-1}110^{t-1}.
\]
A single application of the local rule \eqref{eq:rewrite-adjacent} transforms
this word into $e_{m,t-1}=0^{m-t+1}10^{t-2}\in X_m$, so
Corollary~\ref{cor:fold-order-independent} yields
\[
(\Phi_m(V_m))_t=e_{m,t-1}
\qquad (2\le t\le m-1).
\]
Thus the first $m-1$ output symbols of $\Phi_m(U_m)$ and $\Phi_m(V_m)$ agree.

At time $t=m$, the two raw windows are
\[
010^{m-2}
\qquad\text{and}\qquad
10^{m-1}.
\]
Both words already lie in $X_m$, and they are distinct, proving the displayed
inequality at time $m$.

Now suppose that $\Psi$ has memory $q$ and anticipation $0$ and satisfies
$\Psi\circ\Phi_m=\mathrm{id}$.  Set
\[
\widetilde U_m:=\sigma^{m-1}U_m,
\qquad
\widetilde V_m:=\sigma^{m-1}V_m.
\]
Then
\[
(\widetilde U_m)_0=0,
\qquad
(\widetilde V_m)_0=1,
\]
because $u_m$ and $v_m$ first differ at their $(m-1)$st coordinate.  By
Proposition~\ref{prop:Phi-m-equivariant}, the output sequences
$\Phi_m(\widetilde U_m)$ and $\Phi_m(\widetilde V_m)$ agree on the coordinates
$-(m-2),\ldots,0$.  If $q<m-1$, then the local rule defining $\Psi(\cdot)_0$
reads only a block contained in those common coordinates, so
\[
\Psi(\Phi_m(\widetilde U_m))_0=\Psi(\Phi_m(\widetilde V_m))_0.
\]
Since $\Psi\circ\Phi_m=\mathrm{id}$, this would force
$(\widetilde U_m)_0=(\widetilde V_m)_0$, a contradiction.  Therefore
$q\ge m-1$.
\end{proof}

\begin{table}[t]
\centering
\small
\begin{tabular}{cccc}
\toprule
state $v$ & appended bit $b$ & label $\ell_3(v,b)$ & next state \\
\midrule
$00$ & $0$ & $000$ & $00$ \\
$00$ & $1$ & $001$ & $01$ \\
$01$ & $0$ & $010$ & $10$ \\
$01$ & $1$ & $000$ & $11$ \\
$10$ & $0$ & $100$ & $00$ \\
$10$ & $1$ & $101$ & $01$ \\
$11$ & $0$ & $001$ & $10$ \\
$11$ & $1$ & $100$ & $11$ \\
\bottomrule
\end{tabular}
\caption{The labeled graph $G_3$.}
\label{tab:G3}
\end{table}

\begin{proposition}[Base case at length $3$]
\label{prop:G3-base-case}
Every block
\[
a_1a_2a_3\in\mathcal L_3(Y_3)
\]
has a unique raw lift
\[
\omega_1\cdots\omega_5\in\Omega_5
\]
satisfying
\[
\Fold_3(\omega_j\omega_{j+1}\omega_{j+2})=a_j
\qquad (1\le j\le 3).
\]
Consequently there is a well-defined map
\[
\Theta_{3,3}:\mathcal L_3(Y_3)\to X_1=\{0,1\},
\qquad
\Theta_{3,3}(a_1a_2a_3):=\omega_3.
\]
\end{proposition}

\begin{proof}
Table~\ref{tab:G3} gives the complete labeled graph $G_3$.  A block
$a_1a_2a_3\in\mathcal L_3(Y_3)$ is carried by some length-$3$ path in $G_3$.
Since $G_3$ has four states and two outgoing edges from each state, there are
\[
4\cdot 2^3=32
\]
length-$3$ paths in total.  A direct inspection of the $32$ paths encoded by
Table~\ref{tab:G3} shows that no two distinct length-$3$ paths carry the same
label word.  Hence every admissible block $a_1a_2a_3$ determines a unique
length-$3$ path in $G_3$.

If the initial vertex of that path is $\omega_1\omega_2$ and its successive
appended bits are $\omega_3,\omega_4,\omega_5$, then the path corresponds to
the unique raw lift $\omega_1\cdots\omega_5$.  The central bit $\omega_3$ is
therefore uniquely determined, and this defines $\Theta_{3,3}$.
\end{proof}



\section{Finite-window rigidity and consequences}

We now prove the main reconstruction theorem.  Once one knows that a raw block
is determined by overlapping folded windows, the block counts, local inverse
rules, push-forward measure formulas, and minimal presentation results follow
from the same argument.

\subsection{Reconstruction from overlaps}

Here $\mathcal L_n(Y_m)$ denotes the set of length-$n$ blocks occurring in
$Y_m$.

\begin{lemma}[A signed Fibonacci gap shifts by one index]
\label{lem:signed-fibonacci-gap-shift}
Let $m\ge 2$ and let $d_1,\ldots,d_{m-1}\in\{-1,0,1\}$.  If
\[
\sum_{k=1}^{m-1} d_kF_{k+1}=F_m,
\]
then
\[
\sum_{k=1}^{m-1} d_kF_k=F_{m-1}.
\]
If instead
\[
\sum_{k=1}^{m-1} d_kF_{k+1}=-F_m,
\]
then
\[
\sum_{k=1}^{m-1} d_kF_k=-F_{m-1}.
\]
\end{lemma}

\begin{proof}
It is enough to prove the positive case, since the negative case follows by
replacing each $d_k$ with $-d_k$.

Let
\[
\varphi:=\frac{1+\sqrt5}{2},
\qquad
\psi:=\frac{1-\sqrt5}{2}.
\]
By Binet's formula,
\[
F_n=\frac{\varphi^n-\psi^n}{\sqrt5}
\qquad (n\ge 1).
\]
Define
\[
P(z):=\sum_{k=1}^{m-1} d_kz^k.
\]
The hypothesis reads
\[
\frac{\varphi P(\varphi)-\psi P(\psi)}{\sqrt5}=F_m
=
\frac{\varphi^m-\psi^m}{\sqrt5},
\]
so
\[
C:=\varphi P(\varphi)-\varphi^m
=
\psi P(\psi)-\psi^m.
\]
The first expression is an algebraic integer in $\ZZ[\varphi]$, while the
second is its Galois conjugate.  Since they are equal, $C$ is fixed by
conjugation and therefore lies in $\mathbb{Q}$; hence $C$ is an ordinary
integer.

On the other hand,
\[
|C|
\le
\sum_{k=1}^{m-1} |\psi|^{k+1}+|\psi|^m
<
\sum_{j=2}^{\infty} |\psi|^j
=
1.
\]
Thus the integer $C$ satisfies $|C|<1$, so $C=0$.  Therefore
\[
\varphi P(\varphi)=\varphi^m,
\qquad
\psi P(\psi)=\psi^m,
\]
and hence
\[
P(\varphi)=\varphi^{m-1},
\qquad
P(\psi)=\psi^{m-1}.
\]
Applying Binet's formula again gives
\[
\sum_{k=1}^{m-1} d_kF_k
=
\frac{P(\varphi)-P(\psi)}{\sqrt5}
=
\frac{\varphi^{m-1}-\psi^{m-1}}{\sqrt5}
=
F_{m-1}.
\]
\end{proof}

\begin{proposition}[Pairwise overlap reduction]
\label{prop:pair-overlap-reduction}
For every $m\ge 3$ there is a well-defined map
\[
\Theta_m:\mathcal L_2(Y_m)\to X_{m-1}
\]
with the following property.  If
\[
a=\Fold_m(pr),
\qquad
b=\Fold_m(rq)
\]
for some bits $p,q\in\{0,1\}$ and some overlap word
$r=r_1\cdots r_{m-1}\in\Omega_{m-1}$, then
\[
\Theta_m(ab)=\Fold_{m-1}(r).
\]
\end{proposition}

\begin{proof}
Take two raw words
\[
prq,\qquad p'r'q'
\]
of length $m+1$ such that
\[
\Fold_m(pr)=\Fold_m(p'r'),
\qquad
\Fold_m(rq)=\Fold_m(r'q').
\]
Write
\[
\Delta:=N(r)-N(r').
\]
Since $r,r'\in\Omega_{m-1}$, one has
\[
|\Delta|\le \sum_{k=1}^{m-1}F_{k+1}=F_{m+2}-2<F_{m+2}.
\]
By Corollary~\ref{cor:visible-value-overflow}, equality of the second folded
symbols implies
\[
\Delta+(q-q')F_{m+1}\in\{-F_{m+2},0,F_{m+2}\}.
\]
Because $q-q'\in\{-1,0,1\}$ and $F_{m+2}=F_{m+1}+F_m$, this forces
\[
\Delta\in\{0,\pm F_m,\pm F_{m+1}\}.
\]

We claim that the cases $\Delta=\pm F_m$ are impossible.  Suppose first that
\[
\Delta=F_m.
\]
Set
\[
d_k:=r_k-r_k'\in\{-1,0,1\}
\qquad (1\le k\le m-1).
\]
Then
\[
\sum_{k=1}^{m-1} d_kF_{k+1}=F_m,
\]
so Lemma~\ref{lem:signed-fibonacci-gap-shift} gives
\[
\sum_{k=1}^{m-1} d_kF_k=F_{m-1}.
\]
Hence
\[
\begin{aligned}
N(pr)-N(p'r')
&=
(p-p')+\sum_{k=1}^{m-1} d_kF_{k+2}\\
&=
(p-p')+\sum_{k=1}^{m-1} d_k(F_{k+1}+F_k)\\
&=
(p-p')+F_m+F_{m-1}\\
&=
(p-p')+F_{m+1}.
\end{aligned}
\]
This lies in
\[
\{F_{m+1}-1,F_{m+1},F_{m+1}+1\},
\]
which is disjoint from $\{-F_{m+2},0,F_{m+2}\}$ because
$F_{m+2}=F_{m+1}+F_m>F_{m+1}+1$.  By
Corollary~\ref{cor:visible-value-overflow}, the first folded symbols cannot be
equal, a contradiction.  The case $\Delta=-F_m$ is identical after changing
signs.

Therefore
\[
\Delta\in\{0,\pm F_{m+1}\}.
\]
By Corollary~\ref{cor:visible-value-overflow} at length $m-1$, this is exactly
the condition
\[
\Fold_{m-1}(r)=\Fold_{m-1}(r').
\]
Thus the folded overlap depends only on the pair of output symbols, and the map
$\Theta_m$ is well defined.
\end{proof}

\begin{remark}
The restriction $m\ge 3$ is essential.  For $m=2$, both raw words $00$ and
$11$ fold to $00$, so the pair $(00,00)\in\mathcal L_2(Y_2)$ does not
determine the overlapping bit.
\end{remark}

\begin{definition}[Recursive overlap decoders]
\label{def:recursive-overlap-decoder}
For integers $m\ge 3$ and $1\le n\le m$, define
\[
\Theta_{m,n}:\mathcal L_n(Y_m)\to X_{m-n+1}
\]
as follows.  For $n=1$, set
\[
\Theta_{m,1}(a_1):=a_1.
\]
For $(m,n)=(3,2)$, set
\[
\Theta_{3,2}(a_1a_2):=\Theta_3(a_1a_2).
\]
For $(m,n)=(3,3)$, let $\omega_1\cdots\omega_5$ be the unique raw lift from
Proposition~\ref{prop:G3-base-case} and define
\[
\Theta_{3,3}(a_1a_2a_3):=\omega_3.
\]
For $m\ge 4$ and $2\le n\le m$, define
\[
\Theta_{m,n}(a_1\cdots a_n)
:=
\Theta_{m-1,n-1}\bigl(
\Theta_m(a_1a_2)\,
\Theta_m(a_2a_3)\cdots
\Theta_m(a_{n-1}a_n)
\bigr).
\]
\end{definition}

\begin{proposition}[Common-overlap decoding]
\label{prop:common-overlap-decoder}
Let $m\ge 3$ and $1\le n\le m$.  Suppose
\[
\omega=\omega_1\cdots\omega_{m+n-1}\in\Omega_{m+n-1},
\qquad
a_j:=\Fold_m(\omega_j\cdots\omega_{j+m-1})
\quad (1\le j\le n).
\]
Then
\[
\Theta_{m,n}(a_1\cdots a_n)
=
\Fold_{m-n+1}(\omega_n\cdots\omega_m).
\]
Thus $\Theta_{m,n}$ recovers the fold of the common overlap of the $n$
consecutive raw windows.
\end{proposition}

\begin{proof}
We argue by induction on $m$.

If $m=3$, there are three cases.  For $n=1$, the claim is the definition of
$\Theta_{3,1}$.  For $n=2$, it is exactly
Proposition~\ref{prop:pair-overlap-reduction}.  For $n=3$,
Proposition~\ref{prop:G3-base-case} gives a unique raw lift
\[
\omega_1\cdots\omega_5
\]
of $a_1a_2a_3$, and by definition
\[
\Theta_{3,3}(a_1a_2a_3)=\omega_3.
\]
Since $X_1=\{0,1\}$ and $\Fold_1$ fixes admissible words by
Proposition~\ref{prop:fold-basic}\textup{(ii)}, this is
\[
\Fold_1(\omega_3).
\]
Thus the statement holds for $m=3$.

Now assume $m\ge 4$ and that the proposition is known at level $m-1$ for every
admissible block length.  The case $n=1$ is again the definition.  Assume
$n\ge 2$.  By Proposition~\ref{prop:pair-overlap-reduction}, for each
$1\le j\le n-1$ one has
\[
\Theta_m(a_ja_{j+1})
=
\Fold_{m-1}(\omega_{j+1}\cdots\omega_{j+m-1}).
\]
These are precisely the $(m-1)$-folds of the $n-1$ consecutive length-$(m-1)$
windows cut from the raw block
\[
\omega_2\cdots\omega_{m+n-2}.
\]
Applying the induction hypothesis at level $m-1$ to the $n-1$ folded
length-$(m-1)$ windows yields
\[
\Theta_{m,n}(a_1\cdots a_n)
=
\Theta_{m-1,n-1}\bigl(
\Theta_m(a_1a_2)\cdots\Theta_m(a_{n-1}a_n)
\bigr)
=
\Fold_{m-n+1}(\omega_n\cdots\omega_m),
\]
as claimed.
\end{proof}

\begin{theorem}[Unique lift of a length-$m$ output block]
\label{thm:length-m-block-unique-lift}
Let $m\ge 3$.  Every block
\[
a_1\cdots a_m\in\mathcal L_m(Y_m)
\]
has a unique raw lift
\[
\omega_1\cdots\omega_{2m-1}\in\Omega_{2m-1}
\]
satisfying
\[
\Fold_m(\omega_j\cdots\omega_{j+m-1})=a_j
\qquad (1\le j\le m).
\]
\end{theorem}

\begin{proof}
Existence is immediate from the definition of $\mathcal L_m(Y_m)$.

For uniqueness, suppose
\[
\omega_1\cdots\omega_{2m-1}
\]
is any lift of the given output block.  For $1\le k\le m$ define
\[
L_k:=\Theta_{m,m-k+1}(a_1\cdots a_{m-k+1}),
\qquad
R_k:=\Theta_{m,m-k+1}(a_k\cdots a_m).
\]
By Proposition~\ref{prop:common-overlap-decoder},
\[
L_k=\Fold_k(\omega_{m-k+1}\cdots\omega_m),
\qquad
R_k=\Fold_k(\omega_m\cdots\omega_{m+k-1}).
\]

For $k=1$ this gives
\[
L_1=R_1=\omega_m\in\{0,1\},
\]
so the central bit is determined.

Now fix $2\le k\le m$ and assume inductively that the central blocks
\[
\omega_{m-k+2}\cdots\omega_m,
\qquad
\omega_m\cdots\omega_{m+k-2}
\]
have already been determined.  The word $L_k$ is the fold of the raw $k$-block
\[
\omega_{m-k+1}\omega_{m-k+2}\cdots\omega_m,
\]
whose suffix of length $k-1$ is known.  By
Proposition~\ref{prop:one-step-left-resolving}, applied with $k$ in place of
$m$, the two possible prepended bits $0$ and $1$ produce distinct folds.
Therefore $L_k$ determines the new left bit $\omega_{m-k+1}$ uniquely.

Similarly, $R_k$ is the fold of
\[
\omega_m\cdots\omega_{m+k-2}\omega_{m+k-1},
\]
whose prefix of length $k-1$ is known.  By
Proposition~\ref{prop:one-step-right-resolving}, again with $k$ in place of
$m$, the appended bit is uniquely determined by $R_k$.  Hence $\omega_{m+k-1}$
is fixed as well.

Proceeding outward from $k=1$ to $k=m$ reconstructs all bits
\[
\omega_1,\ldots,\omega_{2m-1}
\]
uniquely.
\end{proof}

\begin{corollary}[Uniform synchronizing length]
\label{cor:uniform-synchronizing-length}
Let $m\ge 3$.  Every block in $\mathcal L_m(Y_m)$ is synchronizing for the
presentation $G_m$.  Moreover the block
\[
e_{m,0}e_{m,1}\cdots e_{m,m-2}
\]
of length $m-1$ is not synchronizing.  Consequently the minimal synchronizing
length of $G_m$ is exactly $m$.
\end{corollary}

\begin{proof}
By Theorem~\ref{thm:length-m-block-unique-lift}, every block in
$\mathcal L_m(Y_m)$ is carried by exactly one length-$m$ path in $G_m$, hence
is synchronizing.

For the second claim, Proposition~\ref{prop:one-sided-inverse-memory-lower-bound}
constructs two sequences $U_m$ and $V_m$ whose first $m-1$ output symbols are
\[
e_{m,0}e_{m,1}\cdots e_{m,m-2}
\]
and coincide there.  The corresponding length-$(m-1)$ paths in $G_m$ end at the
suffix vertices
\[
010^{m-3}
\qquad\text{and}\qquad
10^{m-2},
\]
which are distinct.  Thus the block admits at least two lifts with different
terminal vertices, so it is not synchronizing.  The distinct $m$th output
symbols constructed in Proposition~\ref{prop:one-sided-inverse-memory-lower-bound}
show in particular that these two lifts do not merge after one further step.
\end{proof}

\begin{theorem}[Finite-block losslessness in Fibonacci numeration]
\label{thm:block-bijection-threshold}
Let $m\ge 3$ and $n\ge m$.  Define
\[
B_{m,n}:\Omega_{n+m-1}\to\mathcal L_n(Y_m),
\qquad
B_{m,n}(\omega_1\cdots\omega_{n+m-1}):=a_1\cdots a_n,
\]
where
\[
a_j:=\Fold_m(\omega_j\cdots\omega_{j+m-1})
\qquad (1\le j\le n).
\]
Then $B_{m,n}$ is a bijection.
\end{theorem}

\begin{proof}
Surjectivity is immediate from the definition of $\mathcal L_n(Y_m)$.

For injectivity, fix a block
\[
a_1\cdots a_n\in\mathcal L_n(Y_m).
\]
Its prefix
\[
a_1\cdots a_m\in\mathcal L_m(Y_m)
\]
has a unique raw lift by Theorem~\ref{thm:length-m-block-unique-lift}, hence it
is carried by a unique length-$m$ path in $G_m$.  Since $G_m$ is right-resolving
by Proposition~\ref{prop:one-step-right-resolving}, each subsequent symbol
$a_{m+1},\ldots,a_n$ determines the next edge uniquely.  Thus the whole block
$a_1\cdots a_n$ is carried by a unique length-$n$ path in $G_m$.

Now a length-$n$ path in $G_m$ is uniquely determined by its initial vertex
$v=v_1\cdots v_{m-1}\in V_m$ together with the appended bits
$b_1,\ldots,b_n\in\{0,1\}$.  Equivalently, it is uniquely determined by the raw
word
\[
v_1\cdots v_{m-1}b_1\cdots b_n\in\Omega_{n+m-1}.
\]
By construction of the labels in $G_m$, this raw word is precisely the unique
length-$(n+m-1)$ block whose consecutive folded windows are
$a_1,\ldots,a_n$.  Therefore $B_{m,n}$ is injective.
\end{proof}

\begin{corollary}[Exact block counts above the threshold]
\label{cor:block-counts-threshold}
For every $m\ge 3$ and $n\ge m$,
\[
\lvert\mathcal L_n(Y_m)\rvert=2^{n+m-1}.
\]
In particular,
\[
\lvert\mathcal L_m(Y_m)\rvert=2^{2m-1}.
\]
\end{corollary}

\begin{proof}
Theorem~\ref{thm:block-bijection-threshold} identifies $\mathcal L_n(Y_m)$
with $\Omega_{n+m-1}=\{0,1\}^{n+m-1}$, whose cardinality is $2^{n+m-1}$.
\end{proof}

\begin{table}[t]
\centering
\small
\begin{tabular}{ccc}
\toprule
$(m,n)$ & direct enumeration of $\lvert\mathcal L_n(Y_m)\rvert$ & $2^{n+m-1}$ \\
\midrule
$(3,3)$ & $32$ & $32$ \\
$(3,4)$ & $64$ & $64$ \\
$(3,5)$ & $128$ & $128$ \\
$(4,4)$ & $128$ & $128$ \\
$(4,5)$ & $256$ & $256$ \\
$(4,6)$ & $512$ & $512$ \\
\bottomrule
\end{tabular}
\caption{Small-scale verification of Corollary~\ref{cor:block-counts-threshold}
obtained by brute-force enumeration of all raw blocks in $\Omega_{n+m-1}$ and
folding each consecutive length-$m$ window.}
\label{tab:small-block-counts}
\end{table}

\begin{corollary}[Higher-block identification with the full two-shift]
\label{cor:higher-block-identification}
Let $m\ge 3$ and $n\ge m$.  The bijection $B_{m,n}$ identifies the
$n$-block presentation of $Y_m$ with the $(n+m-1)$-block presentation of the
full binary shift.  Equivalently, consecutive blocks in $\mathcal L_n(Y_m)$
correspond exactly to consecutive overlapping blocks in $\Omega_{n+m-1}$.
\end{corollary}

\begin{proof}
Let
\[
a=a_1\cdots a_n,
\qquad
a'=a_2\cdots a_{n+1}
\]
be consecutive $n$-blocks in some point of $Y_m$.  If
\[
B_{m,n}^{-1}(a)=\omega_1\cdots\omega_{n+m-1},
\qquad
B_{m,n}^{-1}(a')=\omega'_1\cdots\omega'_{n+m-1},
\]
then both raw words are lifts of the same length-$(n+1)$ output block
$a_1\cdots a_{n+1}\in\mathcal L_{n+1}(Y_m)$.  By
Theorem~\ref{thm:block-bijection-threshold}, that length-$(n+1)$ output block
has a unique raw lift in $\Omega_{n+m}$, so necessarily
\[
\omega'_1\cdots\omega'_{n+m-2}=\omega_2\cdots\omega_{n+m-1}.
\]
Thus adjacency of $n$-blocks in $Y_m$ is exactly overlap by $n+m-2$ bits of
their raw lifts in $\Omega_{n+m-1}$, which is the adjacency rule in the
$(n+m-1)$-block presentation of the full binary shift.  Hence the two
higher-block presentations are canonically identified.
\end{proof}

\begin{corollary}[Linear-time finite-block decoding]
\label{cor:linear-time-block-decoding}
Fix $m\ge 3$.  For every $n\ge m$, a block
\[
a_1\cdots a_n\in\mathcal L_n(Y_m)
\]
can be decoded to its unique raw lift in deterministic time $O_m(n)$ and
working space $O(m)$ beyond the space used to store the output.
\end{corollary}

\begin{proof}
By Theorem~\ref{thm:length-m-block-unique-lift}, the initial block
$a_1\cdots a_m$ determines the unique raw lift
\[
\omega_1\cdots\omega_{2m-1}.
\]
This fixes the current suffix vertex
\[
\omega_{m+1}\cdots\omega_{2m-1}\in V_m.
\]
For each $j=m+1,\ldots,n$, the current suffix vertex together with the next
symbol $a_j$ determines a unique outgoing edge of $G_m$ by
Proposition~\ref{prop:one-step-right-resolving}; following that edge reveals
the next appended bit $\omega_{j+m-1}$.  Repeating this step decodes the entire
raw lift.  Since $m$ is fixed, the initialization cost is $O_m(1)$, each new
output symbol is processed in constant time, and only the current
length-$(m-1)$ suffix together with the already decoded raw word must be
stored.
\end{proof}

\begin{remark}[Small decoding examples]
\label{rem:small-decoding-examples}
For $m=3$, the output block
\[
(a_1,a_2,a_3)=(001,010,100)
\]
has unique raw lift $00100$.  Thus the central-bit inverse returns
\[
\kappa_{3,1}(001,010,100)=1.
\]
For $m=4$, the output block
\[
(a_1,a_2,a_3,a_4)=(1001,0010,1000,0000)
\]
has unique raw lift $1110000$.  Hence the position-parameterized inverse family
reads off
\[
\kappa_{4,0}=1,\qquad
\kappa_{4,1}=1,\qquad
\kappa_{4,2}=1,\qquad
\kappa_{4,3}=0
\]
on this single normalized block.  The first example is already fully
admissible, while the second passes through nontrivial local normalization
before decoding.
\end{remark}

\begin{theorem}[Exact finite-type order over the folded alphabet]
\label{thm:Y-m-exact-SFT-order}
For every $m\ge 3$, the shift $Y_m$ is an $m$-step shift of finite type over
the alphabet $X_m$, and it is not an $(m-1)$-step shift of finite type.
\end{theorem}

\begin{proof}
Let
\[
y=(y_t)_{t\in\ZZ}\in X_m^{\ZZ}
\]
be such that every length-$(m+1)$ block of $y$ lies in
$\mathcal L_{m+1}(Y_m)$.  For each $t\in\ZZ$, the length-$m$ block
\[
b_t:=y_{t-m+1}\cdots y_t
\]
belongs to $\mathcal L_m(Y_m)$ and is therefore synchronizing by
Corollary~\ref{cor:uniform-synchronizing-length}.  Let $v_t$ denote the unique
terminal vertex of a path in $G_m$ labeled by $b_t$.

Now fix $t$.  Since
\[
y_{t-m+1}\cdots y_ty_{t+1}\in\mathcal L_{m+1}(Y_m),
\]
there exists a path in $G_m$ carrying this length-$(m+1)$ block.  Its prefix
of length $m$ is labeled by $b_t$, so by uniqueness its terminal vertex is
$v_t$.  Its suffix of length $m$ is labeled by
\[
b_{t+1}=y_{t-m+2}\cdots y_{t+1},
\]
so by the same uniqueness statement its terminal vertex is $v_{t+1}$.  Hence
the last edge of that path is an edge
\[
v_t\xrightarrow{\,y_{t+1}\,}v_{t+1}.
\]
Thus the vertices $(v_t)_{t\in\ZZ}$ form a bi-infinite path in $G_m$ labeled by
$y$, and therefore $y\in Y_m$.  This proves that $Y_m$ is $m$-step.

To show that $Y_m$ is not $(m-1)$-step, let
\[
b:=e_{m,0}e_{m,1}\cdots e_{m,m-2}\in\mathcal L_{m-1}(Y_m).
\]
By Corollary~\ref{cor:uniform-synchronizing-length}, the block $b$ is not
synchronizing, so reading $b$ in $G_m$ can end at at least two distinct
vertices.  By Proposition~\ref{prop:one-sided-inverse-memory-lower-bound}, the
length-$m$ block
\[
bc=b\,010^{m-2}
\]
belongs to $\mathcal L_m(Y_m)$.  By
Theorem~\ref{thm:length-m-block-unique-lift}, this length-$m$ block is carried
by exactly one path in $G_m$.  Consequently, among the terminal vertices
reached by reading $b$, exactly one admits an outgoing edge labeled
$010^{m-2}$.  Since $b$ is not synchronizing, at least one other terminal
vertex reached by $b$ does not admit that label.

Every finite path in $G_m$ occurs in some bi-infinite labeled path because the
underlying graph is the binary de Bruijn graph on words of length $m-1$ and is
therefore strongly connected.  Hence there exist two occurrences of the same
length-$(m-1)$ block $b$ in points of $Y_m$ such that one occurrence may be
followed by $010^{m-2}$ and another may not.  Thus the admissible next symbol
after $b$ depends on more than the last $m-1$ output symbols, so $Y_m$ is not
an $(m-1)$-step shift of finite type.
\end{proof}

\begin{theorem}[Local inverse rules]
\label{thm:Phi-m-conjugacy-threshold}
For every $m\ge 3$ and every $0\le r\le m-1$, let
\[
\kappa_{m,r}:\mathcal L_m(Y_m)\to\{0,1\}
\]
send a block $a_1\cdots a_m$ to the $(r+1)$st bit of its unique raw lift from
Theorem~\ref{thm:length-m-block-unique-lift}.  Define
\[
\Psi_m^{(r)}:Y_m\to\{0,1\}^{\ZZ},
\qquad
(\Psi_m^{(r)}(y))_t:=\kappa_{m,r}(y_{t-r}\cdots y_{t-r+m-1}).
\]
Then $\Psi_m^{(r)}$ is a sliding block code with memory $r$ and anticipation
$m-1-r$, and
\[
\Psi_m^{(r)}\circ\Phi_m=\mathrm{id}_{\{0,1\}^{\ZZ}},
\qquad
\Phi_m\circ\Psi_m^{(r)}=\mathrm{id}_{Y_m}.
\]
Equivalently, each length-$m$ observation window yields an exact inverse rule
for $\Phi_m:\{0,1\}^{\ZZ}\to Y_m$, and $\Phi_m$ is a topological conjugacy
between the full binary shift and the image shift $Y_m$.
\end{theorem}

\begin{proof}
By definition, $(\Psi_m^{(r)}(y))_t$ depends only on the block
$y_{t-r}\cdots y_{t-r+m-1}$, so $\Psi_m^{(r)}$ has memory $r$ and anticipation
$m-1-r$.

Let $s\in\{0,1\}^{\ZZ}$ and set $y:=\Phi_m(s)$.  For each $t\in\ZZ$, applying
Theorem~\ref{thm:length-m-block-unique-lift} to the output block
\[
y_{t-r}\cdots y_{t-r+m-1}
\]
shows that its unique raw lift is
\[
s_{t-r}\cdots s_{t-r+2m-2}.
\]
Therefore the $(r+1)$st bit of that lift is $s_t$, so
\[
(\Psi_m^{(r)}(y))_t=s_t.
\]
Hence
\[
\Psi_m^{(r)}(\Phi_m(s))=s,
\]
so $\Psi_m^{(r)}\circ\Phi_m=\mathrm{id}$.

Now let $y\in Y_m$.  By definition of $Y_m$ there exists
$s\in\{0,1\}^{\ZZ}$ with $\Phi_m(s)=y$.  Applying the first identity to this
$s$ gives
\[
\Psi_m^{(r)}(y)=\Psi_m^{(r)}(\Phi_m(s))=s,
\]
and therefore
\[
\Phi_m(\Psi_m^{(r)}(y))=\Phi_m(s)=y.
\]
Thus $\Phi_m$ and $\Psi_m^{(r)}$ are mutual inverses.
\end{proof}

\begin{corollary}[The one-sided memory threshold is exact]
\label{cor:Phi-m-memory-threshold-exact}
Let $m\ge 3$.  Among all sliding block codes
\[
\Psi:Y_m\to\{0,1\}^{\ZZ}
\]
with anticipation $0$ and
\[
\Psi\circ\Phi_m=\mathrm{id}_{\{0,1\}^{\ZZ}},
\]
the smallest possible memory is exactly $m-1$.
\end{corollary}

\begin{proof}
The endpoint inverse $\Psi_m^{(m-1)}$ from
Theorem~\ref{thm:Phi-m-conjugacy-threshold} has memory $m-1$, while
Proposition~\ref{prop:one-sided-inverse-memory-lower-bound} shows that no
one-sided left inverse can use smaller memory.
\end{proof}

\subsection{Bernoulli measures and block statistics}

We now pass from exact block transport to explicit probability laws.  The block
bijection $B_{m,n}$ identifies each admissible folded block with one raw
binary block of length $n+m-1$, so Bernoulli product measures push forward to
closed cylinder laws on the folded alphabet.

\begin{theorem}[Bernoulli transport and exact cylinder law]
\label{thm:bernoulli-transport-cylinder-law}
Fix $m\ge 3$ and let $\nu_p$ be the Bernoulli product measure on
$\{0,1\}^{\ZZ}$ with parameter $p\in(0,1)$.  Set
\[
\mu_{m,p}:=(\Phi_m)_*\nu_p.
\]
For a finite word $w$ in either shift space, write $[w]_0$ for the cylinder on
the coordinates $0,\ldots,\lvert w\rvert-1$.  Then for every $n\ge m$ and
every $a\in\mathcal L_n(Y_m)$,
\[
\mu_{m,p}([a]_0)
=
\nu_p([\omega(a)]_0)
=
p^{|\omega(a)|_1}(1-p)^{n+m-1-|\omega(a)|_1},
\qquad
\omega(a):=B_{m,n}^{-1}(a),
\]
where $|\omega(a)|_1$ denotes the number of $1$'s in the raw block
$\omega(a)$.

In particular, when $p=\tfrac12$ one has
\[
\mu_{m,1/2}([a]_0)=2^{-(n+m-1)}
\qquad (a\in\mathcal L_n(Y_m),\ n\ge m),
\]
so every admissible folded block of length $n\ge m$ is equally likely.
\end{theorem}

\begin{proof}
By definition of the pushforward measure,
\[
\mu_{m,p}([a]_0)=\nu_p(\Phi_m^{-1}([a]_0)).
\]
Now $\Phi_m(s)\in[a]_0$ means precisely that the raw block
$s_0\cdots s_{n+m-2}\in\Omega_{n+m-1}$ maps to $a$ under $B_{m,n}$.  Since
Theorem~\ref{thm:block-bijection-threshold} makes $B_{m,n}$ bijective, this is
equivalent to
\[
s_0\cdots s_{n+m-2}=\omega(a):=B_{m,n}^{-1}(a).
\]
Hence
\[
\Phi_m^{-1}([a]_0)=[\omega(a)]_0.
\]
The Bernoulli product law of that raw cylinder is
\[
\nu_p([\omega(a)]_0)
=
p^{|\omega(a)|_1}(1-p)^{n+m-1-|\omega(a)|_1},
\]
which proves the formula.  The case $p=\tfrac12$ is immediate.
\end{proof}

\begin{corollary}[Exact entropy profile and excess entropy]
\label{cor:exact-entropy-profile}
Let $(Y_t)_{t\in\ZZ}$ be the coordinate process on $Y_m$ under $\mu_{m,p}$, and
write
\[
h_2(p):=-p\log_2 p-(1-p)\log_2(1-p).
\]
Then for every $n\ge m$,
\[
H_{\mu_{m,p}}(Y_0^{n-1})=(n+m-1)h_2(p).
\]
Consequently
\[
h_{\mu_{m,p}}(\sigma)=h_2(p),
\qquad
E(\mu_{m,p})=(m-1)h_2(p),
\]
where $E(\mu_{m,p})$ denotes the excess entropy.  In particular, for
$p=\tfrac12$ one has
\[
H_{\mu_{m,1/2}}(Y_0^{n-1})=n+m-1,
\qquad
E(\mu_{m,1/2})=m-1.
\]
\end{corollary}

\begin{proof}
Fix $n\ge m$.  By Theorem~\ref{thm:bernoulli-transport-cylinder-law}, the map
$a\mapsto \omega(a)=B_{m,n}^{-1}(a)$ is a bijection from $\mathcal L_n(Y_m)$
onto $\Omega_{n+m-1}$, and the probability of the cylinder $[a]_0$ under
$\mu_{m,p}$ is exactly the Bernoulli probability of the raw cylinder
$[\omega(a)]_0$.  Therefore
\[
H_{\mu_{m,p}}(Y_0^{n-1})
=
-\sum_{a\in\mathcal L_n(Y_m)}\mu_{m,p}([a]_0)\log_2\mu_{m,p}([a]_0)
\]
equals the Shannon entropy of a Bernoulli block of length $n+m-1$, namely
$(n+m-1)h_2(p)$.

Dividing by $n$ and letting $n\to\infty$ gives
$h_{\mu_{m,p}}(\sigma)=h_2(p)$.  The exact linear formula also yields
\[
E(\mu_{m,p})
=
\lim_{n\to\infty}\bigl(H_{\mu_{m,p}}(Y_0^{n-1})-nh_{\mu_{m,p}}(\sigma)\bigr)
=(m-1)h_2(p).
\]
The case $p=\tfrac12$ follows from $h_2(\tfrac12)=1$.
\end{proof}

\begin{theorem}[Exact Markov order under Bernoulli transport]
\label{thm:exact-markov-order-bernoulli}
Fix $m\ge 3$ and $p\in(0,1)$.  For every block
\[
\beta=a_1\cdots a_m\in\mathcal L_m(Y_m),
\]
let $v(\beta)\in V_m$ denote the terminal suffix state of the unique path in
$G_m$ carrying $\beta$.  Then for every $t\in\ZZ$ and every $b\in\{0,1\}$,
\[
\mu_{m,p}\bigl(
Y_{t+1}=\Fold_m(v(\beta)b)\,\big|\,Y_{t-m+1}\cdots Y_t=\beta
\bigr)
=
p^b(1-p)^{1-b}.
\]
In particular, under $\mu_{m,p}$ the coordinate process $(Y_t)_{t\in\ZZ}$ is
an $m$-step Markov process.  It is not an $(m-1)$-step Markov process.
\end{theorem}

\begin{proof}
By shift invariance it suffices to prove the conditional formula at time
$t=m-1$.  Fix
\[
\beta=a_1\cdots a_m\in\mathcal L_m(Y_m),
\]
and let
\[
\omega(\beta)=\omega_1\cdots\omega_{2m-1}:=B_{m,m}^{-1}(\beta).
\]
The terminal suffix state of the unique path carrying $\beta$ is therefore
\[
v(\beta)=\omega_{m+1}\cdots\omega_{2m-1}\in V_m.
\]
By Theorem~\ref{thm:bernoulli-transport-cylinder-law}, the event
$\{Y_0\cdots Y_{m-1}=\beta\}$ pulls back under $\Phi_m$ to the raw cylinder
\[
\{S_0\cdots S_{2m-2}=\omega(\beta)\},
\]
where $(S_t)_{t\in\ZZ}$ is the coordinate process on the Bernoulli shift.
Conditioned on that raw cylinder, the bits
$S_m,\ldots,S_{2m-2}$ are fixed and equal to $v(\beta)$, while the next raw bit
$S_{2m-1}$ remains independent of the past with
\[
\nu_p(S_{2m-1}=b)=p^b(1-p)^{1-b}.
\]
The next folded symbol is exactly
\[
Y_m=\Fold_m(v(\beta)S_{2m-1}),
\]
so
\[
\mu_{m,p}\bigl(
Y_m=\Fold_m(v(\beta)b)\,\big|\,Y_0\cdots Y_{m-1}=\beta
\bigr)
=
p^b(1-p)^{1-b}.
\]
This proves the displayed formula.  Since Proposition~\ref{prop:one-step-right-resolving}
shows that the two labels $\Fold_m(v(\beta)0)$ and $\Fold_m(v(\beta)1)$ are
distinct, the law of the next symbol is determined entirely by the last $m$
output symbols.  Hence $(Y_t)$ is $m$-step Markov under $\mu_{m,p}$.

To see that the Markov order is not smaller, suppose instead that $(Y_t)$ were
an $(m-1)$-step Markov process.  Theorem~\ref{thm:bernoulli-transport-cylinder-law}
shows that every admissible cylinder in $Y_m$ has positive $\mu_{m,p}$-measure,
so $\mu_{m,p}$ has full support on $Y_m$.  For a stationary process with full
support, an $(m-1)$-step Markov law has support equal to the
$(m-1)$-step shift of finite type obtained by declaring a length-$m$ block
admissible exactly when it has positive measure: by repeated conditioning, a
longer block has positive probability if and only if each of its consecutive
length-$m$ subblocks does.  Because every admissible length-$m$ block of $Y_m$
has positive measure, that support would be $Y_m$ itself.  This contradicts
Theorem~\ref{thm:Y-m-exact-SFT-order}, which states that $Y_m$ is not an
$(m-1)$-step shift of finite type.  Therefore the Markov order is exactly $m$.
\end{proof}

\begin{corollary}[Conjugacy invariants]
\label{cor:Y-m-conjugacy-collapse}
For every $m,n\ge 3$, the shifts $Y_m$ and $Y_n$ are topologically conjugate.
In particular, for every $m\ge 3$,
\[
h_{\mathrm{top}}(Y_m)=\log 2,
\qquad
\#\mathrm{Fix}(\sigma^n|_{Y_m})=2^n
\quad (n\ge 1),
\qquad
\zeta_{Y_m}(z)=\frac{1}{1-2z}.
\]
\end{corollary}

\begin{proof}
By Theorem~\ref{thm:Phi-m-conjugacy-threshold}, every $Y_m$ with $m\ge 3$ is
conjugate to the full two-shift.  The conjugacy
\[
Y_m\xrightarrow{\ \Psi_m^{(m-1)}\ }\{0,1\}^{\ZZ}\xrightarrow{\ \Phi_n\ }Y_n
\]
gives the first claim, and the entropy, fixed-point counts, and zeta function
are the corresponding invariants of the full two-shift.
\end{proof}

\subsection{The natural right-resolving presentation}

We now return to the subset construction.  The exact decoding theorem shows
that the determinized presentation has a finite transient ambiguity shell and
that its essential core is already visible in the explicit graph $G_m$.  The
next results quantify the resulting deterministic presentation size on the
natural folded alphabet $X_m$.

\begin{definition}[Deterministic candidate-set presentation]
\label{def:Phi-m-right-resolving}
For a nonempty subset $S\subset V_m$ and a label $a\in X_m$, define
\[
\delta_m(S,a):=
\{v'\in V_m:\exists\, e:v\to v' \text{ with } v\in S \text{ and } \ell_m(e)=a\}.
\]
Starting from the state $V_m$, take all reachable nonempty subsets of $V_m$ and
add labeled edges
\[
S\xrightarrow{a}\delta_m(S,a)
\]
whenever $\delta_m(S,a)\neq \varnothing$.  The resulting finite labeled graph
is denoted $H_m^{\mathrm{det}}$.
\end{definition}

\begin{proposition}[Deterministic right-resolving presentation]
\label{prop:Phi-m-right-resolving}
The graph $H_m^{\mathrm{det}}$ is right-resolving and presents the same labeled
language as $G_m$.  Consequently its edge shift has label image $Y_m$.
\end{proposition}

\begin{proof}
By construction, for a fixed state $S$ and label $a$, there is at most one
outgoing edge with label $a$, so $H_m^{\mathrm{det}}$ is right-resolving.  The
subset construction preserves the accepted labeled language of a finite
automaton, hence $H_m^{\mathrm{det}}$ and $G_m$ present the same set of finite
label words.  Their bi-infinite label shifts therefore coincide, and by
Theorem~\ref{thm:Phi-m-sofic-graph} this common image is $Y_m$.
\end{proof}

\begin{proposition}[The singleton core]
\label{prop:deterministic-singleton-core}
Let $m\ge 3$.  Then:
\begin{enumerate}[label=\textup{(\roman*)}]
\item every state of $H_m^{\mathrm{det}}$ reached after at least $m$ input
symbols is a singleton;
\item the reachable singleton states are exactly the sets $\{v\}$ with
$v\in V_m$;
\item the labeled subgraph induced by the singleton states is canonically
labeled-isomorphic to $G_m$.
\end{enumerate}
\end{proposition}

\begin{proof}
Take any block
\[
a_1\cdots a_m\in\mathcal L_m(Y_m).
\]
By Theorem~\ref{thm:length-m-block-unique-lift}, this block has a unique raw
lift, hence a unique length-$m$ path in $G_m$ carrying that label word.  Its
terminal vertex is therefore uniquely determined by the block.  Consequently,
every reachable state after reading $m$ symbols is a singleton.  Since
singletons are forward closed under the subset construction, every state reached
after at least $m$ symbols is again a singleton.  This proves
\textup{(i)}.

For \textup{(ii)}, it remains to show that every singleton $\{v\}$ with
$v\in V_m$ is reachable.  By
Theorem~\ref{thm:alternating-single-symbol-synchronizing}, the symbol $b_m$
sends the full candidate set to a singleton.  The
underlying directed graph of $G_m$ is the binary de Bruijn graph on words of
length $m-1$, so it is strongly connected.  Thus for every $v\in V_m$ there
exists a path in $G_m$ from that singleton state to $v$, and reading its label
from the singleton yields $\{v\}$.

Finally, \textup{(iii)} is immediate from the definitions: for singleton states,
\[
\delta_m(\{v\},a)=\{v'\}
\]
if and only if $G_m$ has an edge $v\to v'$ labeled $a$.  Hence the singleton
subgraph is exactly a relabeling of $G_m$.
\end{proof}

\begin{proposition}[The ambiguity shell is transient]
\label{prop:deterministic-ambiguity-transient}
Let $m\ge 3$.  Every non-singleton state of $H_m^{\mathrm{det}}$ is transient.
Equivalently, the subgraph induced by the non-singleton states is acyclic.
\end{proposition}

\begin{proof}
Suppose there were a directed cycle consisting entirely of non-singleton states.
Since every state of $H_m^{\mathrm{det}}$ is reachable from the initial state
$V_m$, one could first reach that cycle and then traverse it often enough to
obtain a path of length at least $m$ ending at a non-singleton state.  This
contradicts Proposition~\ref{prop:deterministic-singleton-core}\textup{(i)}.
Hence no such cycle exists, and every non-singleton state is transient.
\end{proof}

\begin{proposition}[Exact nilpotency of the ambiguity shell]
\label{prop:ambiguity-shell-nilpotency}
Let $m\ge 3$.  Order the states of $H_m^{\mathrm{det}}$ by placing the
non-singleton states first and the singleton core last.  Then the adjacency
matrix of $H_m^{\mathrm{det}}$ has the form
\[
A(H_m^{\mathrm{det}})
=
\begin{pmatrix}
N_m & R_m\\
0 & A(G_m)
\end{pmatrix}
\]
for some matrix $R_m$, and
\[
N_m^m=0,
\qquad
N_m^{m-1}\neq 0.
\]
Consequently
\[
\det(I-zA(H_m^{\mathrm{det}}))=\det(I-zA(G_m)).
\]
\end{proposition}

\begin{proof}
By Proposition~\ref{prop:deterministic-singleton-core}\textup{(iii)}, the
singleton states induce a labeled copy of $G_m$.  Since singleton states are
forward closed, there are no edges from the singleton core back into the
non-singleton shell, so the adjacency matrix has the displayed block upper
triangular form.

By Proposition~\ref{prop:deterministic-singleton-core}\textup{(i)}, every path
of length $m$ in $H_m^{\mathrm{det}}$ ends in a singleton.  Hence there is no
path of length $m$ contained entirely in the non-singleton shell, which is
equivalent to $N_m^m=0$.

By Corollary~\ref{cor:uniform-synchronizing-length}, there exists a
non-synchronizing block of length $m-1$.  Reading that block from the initial
state $V_m$ yields a path of length $m-1$ ending at a non-singleton state.
Because singleton states are forward closed, every intermediate state on that
path is also non-singleton.  Therefore the shell contains a path of length
$m-1$, so $N_m^{m-1}\neq 0$.

Finally,
\[
\det(I-zA(H_m^{\mathrm{det}}))
=
\det(I-zN_m)\det(I-zA(G_m)).
\]
Since $N_m$ is nilpotent, $\det(I-zN_m)=1$, and the determinant identity
follows.
\end{proof}

\begin{proposition}[The explicit graph is follower separated]
\label{prop:G-m-follower-separated}
Let $m\ge 3$.  Distinct vertices of $G_m$ have distinct follower sets.
\end{proposition}

\begin{proof}
Let
\[
b_m=u_1\cdots u_m
\]
be the alternating symbol from
Definition~\ref{def:alternating-extremal-words}.  For each vertex
\[
v=v_1\cdots v_{m-1}\in V_m,
\]
consider the length-$m$ path obtained by appending the fixed bit word
$u_1\cdots u_m$.  Let
\[
\Lambda(v)\in\mathcal L_m(Y_m)
\]
denote its label word.

We claim that $\Lambda(v)$ is carried by exactly one path in $G_m$, namely the
one that starts at $v$ and appends $u_1,\ldots,u_m$.  Indeed, the last symbol
of $\Lambda(v)$ is $b_m$, and by
Theorem~\ref{thm:alternating-single-symbol-synchronizing}, the symbol $b_m$ is
carried by a unique edge.  Thus the final edge of any path carrying
$\Lambda(v)$ is fixed.  Since $G_m$ is one-step left-resolving by
Proposition~\ref{prop:one-step-left-resolving}, the penultimate label then
determines the penultimate edge uniquely; iterating this backward reconstruction
shows that the whole path carrying $\Lambda(v)$ is unique.

In particular, its initial vertex is uniquely determined by $\Lambda(v)$.
Hence if $v\neq w$, the word $\Lambda(v)$ belongs to the follower set of $v$
but not to that of $w$.  Therefore distinct vertices have distinct follower
sets.
\end{proof}

\begin{corollary}[The graph $G_m$ is the canonical minimal right-resolving presentation]
\label{cor:G-m-fischer-cover}
For every $m\ge 3$, the canonical minimal right-resolving presentation of
$Y_m$ is labeled-isomorphic to $G_m$.  Equivalently, $G_m$ is the right
Fischer cover of $Y_m$.
\end{corollary}

\begin{proof}
By Theorem~\ref{thm:Phi-m-conjugacy-threshold}, the shift $Y_m$ is irreducible
(indeed conjugate to the full two-shift).  Proposition~\ref{prop:Phi-m-right-resolving}
gives a deterministic right-resolving presentation of $Y_m$.  By
Proposition~\ref{prop:deterministic-ambiguity-transient}, the only essential
component of that presentation is the singleton core from
Proposition~\ref{prop:deterministic-singleton-core}, which is labeled-isomorphic
to $G_m$.  Proposition~\ref{prop:G-m-follower-separated} shows that this core is
already follower separated.  The standard construction of the canonical
minimal right-resolving presentation, equivalently of the right Fischer cover,
as the follower-separated essential component of a deterministic
right-resolving presentation therefore leaves $G_m$ unchanged; see
\cite[Chs.~3--4]{LindMarcus1995} and \cite[Ch.~4]{Kitchens1998SymbolicDynamics}.
\end{proof}

\begin{theorem}[Size of the right-resolving presentation]
\label{thm:natural-state-complexity-gap}
For every $m\ge 3$, the canonical minimal right-resolving presentation of
$Y_m$ over the folded alphabet $X_m$ has exactly
\[
\lvert V(G_m)\rvert=2^{m-1},
\qquad
\lvert E(G_m)\rvert=2^m
\]
states and edges.  In particular, no right-resolving presentation of $Y_m$
over $X_m$ uses fewer than $2^{m-1}$ states.  Since
\[
\lvert X_m\rvert=F_{m+2},
\]
one has
\[
\frac{\lvert V(G_m)\rvert}{\lvert X_m\rvert}
=
\frac{2^{m-1}}{F_{m+2}}
\sim
\frac{\sqrt5}{2\varphi^2}\Bigl(\frac{2}{\varphi}\Bigr)^m,
\qquad
\varphi=\frac{1+\sqrt5}{2}.
\]
Thus, although every $Y_m$ with $m\ge 3$ is conjugate to the full two-shift,
its canonical deterministic presentation over the natural folded alphabet has
exponentially growing state complexity, even relative to the alphabet size.
\end{theorem}

\begin{proof}
Corollary~\ref{cor:G-m-fischer-cover} identifies $G_m$ as the canonical minimal
right-resolving presentation of $Y_m$.  The state and edge counts are therefore
those of $G_m$, and Definition~\ref{def:G-m} gives
\[
\lvert V(G_m)\rvert=2^{m-1},
\qquad
\lvert E(G_m)\rvert=2^m.
\]
Minimality of the Fischer cover yields the state lower bound for all
right-resolving presentations over $X_m$.

The alphabet size is
\[
\lvert X_m\rvert=F_{m+2}
\]
by Corollary~\ref{cor:stable-syntax-counting}.  Using the Binet asymptotic
$F_{m+2}\sim \varphi^{m+2}/\sqrt5$ gives
\[
\frac{2^{m-1}}{F_{m+2}}
\sim
\frac{2^{m-1}\sqrt5}{\varphi^{m+2}}
=
\frac{\sqrt5}{2\varphi^2}\Bigl(\frac{2}{\varphi}\Bigr)^m,
\]
which proves the claimed representation gap.
\end{proof}

\begin{corollary}[Coincidence of decoding and presentation scales]
\label{cor:local-complexity-scale-coincidence}
For every $m\ge 3$, the following quantities are exact and all equal to the
same scale:
\begin{enumerate}[label=\textup{(\roman*)}]
\item the finite-type order of $Y_m$ over the alphabet $X_m$ is $m$;
\item the minimal synchronizing length of the presentation $G_m$ is $m$;
\item the nilpotency index of the ambiguity-shell matrix $N_m$ is $m$;
\item the minimal memory of a one-sided sliding-block inverse of $\Phi_m$ is
$m-1$.
\end{enumerate}
In particular, although the shifts $Y_m$ with $m\ge 3$ are all conjugate to the
full two-shift, their natural presentations over the folded alphabets $X_m$
still reflect the window length.  Theorem~\ref{thm:exact-markov-order-bernoulli}
shows that the same threshold is also the Markov memory under every transported
Bernoulli law, while
Theorem~\ref{thm:natural-state-complexity-gap} quantifies the accompanying
exponential presentation gap on the natural alphabet.
\end{corollary}

\begin{proof}
Part \textup{(i)} is Theorem~\ref{thm:Y-m-exact-SFT-order},
part \textup{(ii)} is Corollary~\ref{cor:uniform-synchronizing-length},
part \textup{(iii)} is Proposition~\ref{prop:ambiguity-shell-nilpotency}, and
part \textup{(iv)} is Corollary~\ref{cor:Phi-m-memory-threshold-exact}.
\end{proof}



\section{Arithmetic reformulations and consequences}

The results of the preceding sections admit purely number-theoretic
formulations that do not refer to shift spaces, factor maps, or finite-state
presentations.  Throughout this section fix $m\ge 3$, $n\ge m$, and write
$L:=n+m-1$.

\subsection{Local Fibonacci congruence rigidity}

For a binary block $\omega=\omega_1\cdots\omega_L\in\{0,1\}^L$ and
$1\le t\le n$, define the $t$-th local Fibonacci weighted sum
\[
S_t^{(m)}(\omega):=\sum_{j=0}^{m-1}\omega_{t+j}F_{j+2},
\]
and its standard residue modulo $F_{m+2}$,
\[
R_t^{(m)}(\omega)\in\{0,1,\ldots,F_{m+2}-1\},
\qquad
R_t^{(m)}(\omega)\equiv S_t^{(m)}(\omega)\pmod{F_{m+2}}.
\]
Collect all local residues into a vector
\[
\mathcal R_{m,n}(\omega):=
\bigl(R_1^{(m)}(\omega),\ldots,R_n^{(m)}(\omega)\bigr)
\in\{0,1,\ldots,F_{m+2}-1\}^n.
\]

The key observation is that by Corollary~\ref{cor:visible-value-overflow},
\[
R_t^{(m)}(\omega)=V_m\bigl(\Fold_m(\omega_t\cdots\omega_{t+m-1})\bigr),
\]
where $V_m:X_m\to\{0,1,\ldots,F_{m+2}-1\}$ is the bijection from
Proposition~\ref{prop:visible-value-parametrization}.

\begin{theorem}[Local Fibonacci congruence rigidity]
\label{thm:congruence-rigidity}
The map
$\mathcal R_{m,n}:\{0,1\}^L\to\{0,1,\ldots,F_{m+2}-1\}^n$
is injective.  Equivalently, for any
$(r_1,\ldots,r_n)\in\{0,1,\ldots,F_{m+2}-1\}^n$,
the congruence system
\[
\sum_{j=0}^{m-1}\omega_{t+j}F_{j+2}\equiv r_t\pmod{F_{m+2}}
\qquad (1\le t\le n)
\]
has at most one solution $\omega\in\{0,1\}^L$.
The image $\mathcal R_{m,n}(\{0,1\}^L)$ has exactly $2^L=2^{n+m-1}$ elements.
\end{theorem}

\begin{proof}
Take any $\omega\in\{0,1\}^L$.  For each $t$, set
$a_t:=\Fold_m(\omega_t\cdots\omega_{t+m-1})\in X_m$.  By
Corollary~\ref{cor:visible-value-overflow},
$V_m(a_t)\equiv S_t^{(m)}(\omega)\pmod{F_{m+2}}$, and since $V_m(a_t)$
already lies in $\{0,1,\ldots,F_{m+2}-1\}$, one has
$V_m(a_t)=R_t^{(m)}(\omega)$.  Thus
\[
\mathcal R_{m,n}(\omega)=\bigl(V_m(a_1),\ldots,V_m(a_n)\bigr).
\]

Suppose $\mathcal R_{m,n}(\omega)=\mathcal R_{m,n}(\omega')$.  Then
$V_m(a_t)=V_m(a_t')$ for each $t$.  Since $V_m$ is a bijection
(Proposition~\ref{prop:visible-value-parametrization}), $a_t=a_t'$ for every
$1\le t\le n$.  Hence $\omega$ and $\omega'$ produce the same length-$n$
folded block under $B_{m,n}$.  By
Theorem~\ref{thm:block-bijection-threshold}, $B_{m,n}$ is a bijection, so
$\omega=\omega'$.  The image has $2^L$ elements because $\mathcal R_{m,n}$ is
an injection from a set of size $2^L$.
\end{proof}

\subsection{Fingerprint separation of equi-valued representations}

\begin{corollary}[Fingerprint separation]
\label{cor:fingerprint-separation}
For each integer $N$, let
\[
\Omega_L(N):=
\Bigl\{
\omega\in\{0,1\}^L:
\sum_{k=1}^L\omega_kF_{k+1}=N
\Bigr\}.
\]
Then $\mathcal R_{m,n}$ restricted to $\Omega_L(N)$ is still injective:
\[
\lvert\Omega_L(N)\rvert
=
\bigl\lvert\{\mathcal R_{m,n}(\omega):\omega\in\Omega_L(N)\}\bigr\rvert.
\]
\end{corollary}

\begin{proof}
Since $\mathcal R_{m,n}$ is injective on all of $\{0,1\}^L$
(Theorem~\ref{thm:congruence-rigidity}), it is injective on every subset.
\end{proof}

In other words, all length-$L$ binary Fibonacci representations of a single
integer $N$ have mutually distinct overlapping local Zeckendorf fingerprints.

\subsection{Exact bit recovery from residue windows}

Write
$\mathfrak R_{m,\ell}:=\mathcal R_{m,\ell}(\{0,1\}^{\ell+m-1})$
for the set of realizable local residue vectors of length $\ell$.

\begin{theorem}[Bit recovery from local congruence data]
\label{thm:bit-recovery-residues}
For each $0\le s\le m-1$, there exists a function
$\rho_{m,s}:\mathfrak R_{m,m}\to\{0,1\}$ such that for every
$\omega\in\{0,1\}^{n+m-1}$ and every $1\le t\le n-m+1$,
\[
\omega_{t+s}
=
\rho_{m,s}\bigl(
R_t^{(m)}(\omega),R_{t+1}^{(m)}(\omega),\ldots,R_{t+m-1}^{(m)}(\omega)
\bigr).
\]
In particular, taking $s=m-1$ gives a one-sided decoding rule
$g_m:=\rho_{m,m-1}$ with
\[
\omega_t=g_m\bigl(R_{t-m+1}^{(m)}(\omega),\ldots,R_t^{(m)}(\omega)\bigr).
\]
No rule using fewer than $m$ consecutive residues can recover all input bits.
\end{theorem}

\begin{proof}
For any admissible $m$-tuple $(r_1,\ldots,r_m)\in\mathfrak R_{m,m}$, set
$a_i:=V_m^{-1}(r_i)\in X_m$ and define
\[
\rho_{m,s}(r_1,\ldots,r_m):=\kappa_{m,s}(a_1\cdots a_m),
\]
where $\kappa_{m,s}$ is the local inverse rule from
Theorem~\ref{thm:Phi-m-conjugacy-threshold}.

Now take any $\omega\in\{0,1\}^{n+m-1}$ and fix $t$.  Set
$a_j:=\Fold_m(\omega_j\cdots\omega_{j+m-1})$ for $t\le j\le t+m-1$.  Then
$R_j^{(m)}(\omega)=V_m(a_j)$, so $a_j=V_m^{-1}(R_j^{(m)}(\omega))$.  By
Theorem~\ref{thm:Phi-m-conjugacy-threshold}, the folded block
$a_t\cdots a_{t+m-1}$ has unique raw lift
$\omega_t\cdots\omega_{t+2m-2}$, and $\kappa_{m,s}$ reads its $(s+1)$-st
bit.  Therefore
\[
\rho_{m,s}\bigl(R_t^{(m)}(\omega),\ldots,R_{t+m-1}^{(m)}(\omega)\bigr)
=
\kappa_{m,s}(a_t\cdots a_{t+m-1})
=
\omega_{t+s}.
\]

For the optimality claim, suppose a rule using only $k+1<m$ consecutive
residues existed.  Since $V_m$ is a symbol-by-symbol bijection, this would
yield a sliding-block left inverse $\Psi:Y_m\to\{0,1\}^{\ZZ}$ with memory
$k<m-1$, contradicting
Corollary~\ref{cor:Phi-m-memory-threshold-exact}.
\end{proof}

\subsection{Distribution of random local residue vectors}

\begin{theorem}[Exact distribution of random local residues]
\label{thm:random-residue-distribution}
Let $X_1,\ldots,X_L$ be independent with
$\mathbb P(X_i=1)=p$ and $\mathbb P(X_i=0)=1-p$,
where $0<p<1$.  Define the random local residue vector
$\mathbf R_{m,n}:=\mathcal R_{m,n}(X_1\cdots X_L)$.  Then for every
$\mathbf r=(r_1,\ldots,r_n)\in\{0,1,\ldots,F_{m+2}-1\}^n$,
\[
\mathbb P(\mathbf R_{m,n}=\mathbf r)=
\begin{cases}
p^{|\omega(\mathbf r)|_1}(1-p)^{L-|\omega(\mathbf r)|_1},
& \mathbf r\in\mathfrak R_{m,n},\\[4pt]
0,& \mathbf r\notin\mathfrak R_{m,n},
\end{cases}
\]
where $\omega(\mathbf r)$ is the unique preimage guaranteed by
Theorem~\textup{\ref{thm:congruence-rigidity}}.  When $p=\tfrac12$,
$\mathbf R_{m,n}$ is uniformly distributed on $\mathfrak R_{m,n}$:
\[
\mathbb P(\mathbf R_{m,n}=\mathbf r)=2^{-L}
\qquad (\mathbf r\in\mathfrak R_{m,n}).
\]
\end{theorem}

\begin{proof}
If $\mathbf r\notin\mathfrak R_{m,n}$, then no
$\omega\in\{0,1\}^L$ satisfies $\mathcal R_{m,n}(\omega)=\mathbf r$, so the
probability is $0$.

Suppose $\mathbf r\in\mathfrak R_{m,n}$.  Set
$a_t:=V_m^{-1}(r_t)\in X_m$ for $1\le t\le n$.  The event
$\{\mathbf R_{m,n}=\mathbf r\}$ is equivalent to the folded block being
$a_1\cdots a_n$.  By Theorem~\ref{thm:block-bijection-threshold}, this
folded block has a unique raw lift $\omega(\mathbf r)\in\{0,1\}^L$.  Hence
\[
\{\mathbf R_{m,n}=\mathbf r\}
=
\{(X_1,\ldots,X_L)=\omega(\mathbf r)\}.
\]
Since $X_1,\ldots,X_L$ are independent Bernoulli$(p)$,
\[
\mathbb P(\mathbf R_{m,n}=\mathbf r)
=
p^{|\omega(\mathbf r)|_1}(1-p)^{L-|\omega(\mathbf r)|_1}.
\]
When $p=\tfrac12$, every $\omega\in\{0,1\}^L$ has probability $2^{-L}$.
\end{proof}

\begin{corollary}[Entropy of random local residues]
\label{cor:entropy-random-residues}
Under the assumptions of Theorem~\textup{\ref{thm:random-residue-distribution}},
for every $n\ge m$,
\[
H(\mathbf R_{m,n})=(n+m-1)h_2(p),
\]
where $h_2(p):=-p\log_2 p-(1-p)\log_2(1-p)$.  In particular, when
$p=\tfrac12$, $H(\mathbf R_{m,n})=n+m-1$.
\end{corollary}

\begin{proof}
By Theorem~\ref{thm:random-residue-distribution}, the distribution of
$\mathbf R_{m,n}$ is the image of a length-$L$ Bernoulli block under the
bijection $\mathcal R_{m,n}$.  Entropy is invariant under bijective relabeling,
so $H(\mathbf R_{m,n})=H(X_1,\ldots,X_L)=Lh_2(p)=(n+m-1)h_2(p)$.
\end{proof}

\subsection{Markov order of the residue process}

\begin{theorem}[Exact Markov order of the local residue process]
\label{thm:residue-markov-order}
Let $(X_t)_{t\in\ZZ}$ be a two-sided Bernoulli$(p)$ process with $0<p<1$.
Define
\[
R_t:=
\biggl(\sum_{j=0}^{m-1}X_{t+j}F_{j+2}\biggr)\bmod F_{m+2},
\qquad
R_t\in\{0,1,\ldots,F_{m+2}-1\}.
\]
Then $(R_t)_{t\in\ZZ}$ is an exact $m$-step Markov process: given any
realizable block $(r_{t-m+1},\ldots,r_t)$, the conditional law of $R_{t+1}$
depends only on these $m$ values, and no $(m-1)$-step Markov law can describe
$(R_t)$.

More precisely, let $a_i:=V_m^{-1}(r_{t-m+i})\in X_m$,
$\beta:=a_1\cdots a_m\in\mathcal L_m(Y_m)$, and let
$v(\beta)\in\{0,1\}^{m-1}$ be the terminal suffix of the unique raw lift of
$\beta$.  Then for each $b\in\{0,1\}$,
\[
\mathbb P\!\left(
R_{t+1}=V_m(\Fold_m(v(\beta)b))
\;\middle\vert\;
R_{t-m+1}=r_{t-m+1},\ldots,R_t=r_t
\right)
=
p^b(1-p)^{1-b}.
\]
\end{theorem}

\begin{proof}
Set $Y_t:=\Fold_m(X_t\cdots X_{t+m-1})\in X_m$.  By
Corollary~\ref{cor:visible-value-overflow} and
Proposition~\ref{prop:visible-value-parametrization}, $R_t=V_m(Y_t)$ for all $t$.
Since $V_m$ is a bijection, $(R_t)$ and $(Y_t)$ are related by a
symbol-by-symbol relabeling.

Theorem~\ref{thm:exact-markov-order-bernoulli} states that $(Y_t)$ is an exact
$m$-step Markov process under $\mu_{m,p}$, with the displayed transition
probabilities.  Translating through $V_m$ yields the same conclusion for
$(R_t)$: the event $\{Y_j=a_{j-(t-m)}\}$ is equivalent to $\{R_j=r_j\}$
because $R_j=V_m(Y_j)$ and $V_m$ is invertible.  Likewise
$Y_{t+1}=\Fold_m(v(\beta)b)$ if and only if
$R_{t+1}=V_m(\Fold_m(v(\beta)b))$.

If $(R_t)$ were $(m-1)$-step Markov, then $(Y_t)=V_m^{-1}(R_t)$ would also be
$(m-1)$-step Markov, contradicting
Theorem~\ref{thm:exact-markov-order-bernoulli}.
\end{proof}



\section*{Competing interests}

The author(s) declare none.

\section*{Declaration of generative AI and AI-assisted technologies in the manuscript preparation process}

During the preparation of this work the author(s) used OpenAI ChatGPT and
Anthropic Claude in order to assist with language polishing and editing of the
manuscript text.  After using these tools, the author(s) reviewed and edited
the content as needed and take(s) full responsibility for the content of the
published article.

\bibliographystyle{plain}
\bibliography{../../2026_golden_ratio_driven_scan_projection_generation_recursive_emergence/references}

\end{document}
